% concatenated from polyharmonic_map-version_2.tex

\documentclass[12pt,UTF8]{article}
\usepackage{amssymb,latexsym}
\usepackage{amsmath,amsthm}
\usepackage{indentfirst}
%\usepackage{ctex}
%\usepackage[%% Biblatex >>>2
%style = numeric,
%language = brazil,
%giveninits,
%uniquename = init,
%hyperref,
%labelnumber,            % These two options are used just to show the
%defernumbers = true,    % numbers properly in the abnt-numeric examples
%]{biblatex}
%\addbibresource{E:/myonlybib/myonlymathscinetbibfrom2023.bib}
%\addbibresource{E:/myonlybib/low-quality-bib-to-publish.bib}
%\addbibresource{E:/myonlybib/mybib2023.bib}
\usepackage[colorlinks,backref=pages]{hyperref}


\newcommand{\piant}{e^{-\phi}_{*}\frac{\partial }{\partial t }}
\newcommand{\pians}{e^{-\phi}_{*}\frac{\partial }{\partial s }}
\newcommand{\pian}[2]{\frac{\partial #1}{\partial #2}}
\newcommand{\yuankuo}[1]{\left( #1\right) }
\newcommand{\jiankuo}[1]{\left\langle  #1\right\rangle }
\newcommand{\fanshu}[1]{\left\|   #1\right\right\|  }
\newcommand{\covt}[1]{\tilde{\nabla}_{#1}}
\newcommand{\F}{F^{\prime}\yuankuo{\frac{|du|^2}{2}}}
\newcommand{\FF}{F^{\prime\prime}\yuankuo{\frac{|du|^2}{2}}}
\newcommand{\pl}[1]{e^{-\phi}_*e_{#1}}
\newcommand{\cov}[2]{e^{-\phi}_{*}\nabla_{e_{#1}}  e_{#2}  }
\newcommand{\curtensor}[4]{R^N\left(#1,#2,#3, #4\right) }
\newcommand{\curv}[4]{R_{#1,#2,#3, #4} }
\newcommand{\curva}[4]{R^{#1}_{#2,#3, #4} }
\newcommand{\suanzi}[1]{\operatorname{#1}}
\newcommand{\uu}[2]{u_{#1}^{#2}}
\newcommand{\hessian}[2]{\nabla ^2 u (#1,#2)}
\newcommand{\he}[2]{\sum_{#1}^{#2}}
\newcommand{\jifen}[1]{\int_M #1 \dv}
\newcommand{\h}[2]{h(du(#1),du(#2)) }
%%%%%%%%%%%%%%%%%%%%%%%%%%%%%
%\usepackage[notref,notcite]{showkeys}
%\usepackage{citesort}
%
%\emph{}
\usepackage{color}

\newcommand{\cc}[1]{G^{#1}\left( \frac{\|u^*h\|^2}{4} \right)  }
%\usepackage{txfonts}
%\usepackage{pdfsync}
%\usepackage[colorlinks]{hyperref}
%\usepackage[pdftex,linkcolor=blue,citecolor=blue]{hyperref}
%\setlength{\parindent}{1em}
%\renewcommand{\baselinestretch}{1.2}
%\usepackage{txfonts}
%\usepackage{pdfsync}
%\usepackage{showkeys}
%\usepackage{ctex}
%\usepackage{natbib}
%\usepackage[square, comma]{natbib}
\theoremstyle{plain}
\newtheorem{thm}{Theorem}[section]
%\newtheorem{rj}{Theorem}[section]
\newtheorem{cor}{Corollary}
\newtheorem{lem}{Lemma}
\newtheorem{prop}{Proposition}
\theoremstyle{definition}
\newtheorem{defn}{Definition}
\newtheorem{con}{Conjecture}
\newtheorem{rem}{Remark}

\allowdisplaybreaks
\def\d {\mathrm{d}}
\def\Ric {\mathrm{Ric}}
\def\vol{\mathrm{Vol}}
\def\div{\mathrm{div}}
\def\hess{\mathrm{Hess}}
\def\S{\mathbb{S}}
\def\S{\mathbb{R}}
\def\o{\mathcal{O}}

\def\na{\ensuremath{\nabla}}
\def\N{{\Bbb N}}
\def\Z{{\Bbb Z}}
\def\R{{\Bbb R}}
\def\C{{\Bbb C}}
\def\lcm{\operatorname{lcm}}
\def\Res{\operatorname{Res}}


%\def\equationautorefname{式}%
%\def\footnoteautorefname{脚注}%
%\def\itemautorefname{项}%
%\def\figureautorefname{图}%
%\def\tableautorefname{表}%
%\def\partautorefname{篇}%
%\def\appendixautorefname{附录}%
%\def\chapterautorefname{章}%
%\def\sectionautorefname{节}%
%\def\subsectionautorefname{小小节}%
%\def\subsubsectionautorefname{subsubsection}%
%\def\paragraphautorefname{段落}%
%\def\subparagraphautorefname{子段落}%
%\def\FancyVerbLineautorefname{行}%


\def\thmau{Theorem}%
\def\thmautorefname{Theorem}%
\def\lemautorefname{Lemma}%
\usepackage{titlesec}
%\titleformat{\section}[runin]
%{\normalfont\large\bfseries}{\thesection}{1em}{}
%\titleformat{\subsection}[runin]
%{\normalfont\large\bfseries}{\thesubsection}{1em}{}
\renewcommand{\b}[1]{\mathbf{#1}}
\renewcommand{\c}[1]{\mathcal{#1}}
\renewcommand{\d}[1]{\mathbb{#1}}
\newcommand{\f}[1]{\mathfrak{#1}}
\renewcommand{\r}[1]{\mathrm{#1}}
\newcommand{\s}[1]{\mathscr{#1}}
\renewcommand{\sf}[1]{\mathsf{#1}}
\renewcommand{\(}{\left(}
\renewcommand{\)}{\right)}
\newcommand{\res}{\mathbin{|}}
\newcommand{\ol}[1]{\overline{#1}{}}
\newcommand{\wt}[1]{\widetilde{#1}{}}
\newcommand{\ul}{\underline}
\renewcommand{\leq}{\leqslant}
\renewcommand{\geq}{\geqslant}


%%%%%%%%%%%%%%%%%%%%%%%%%%%%%%%%%%%%%%%%%%%%%

\newcommand{\bA}{\b A}
\newcommand{\bB}{\b B}
\newcommand{\bC}{\b C}
\newcommand{\bD}{\b D}
\newcommand{\bE}{\b E}
\newcommand{\bF}{\b F}
\newcommand{\bG}{\b G}
\newcommand{\bH}{\b H}
\newcommand{\bI}{\b I}
\newcommand{\bJ}{\b J}
\newcommand{\bK}{\b K}
\newcommand{\bL}{\b L}
\newcommand{\bM}{\b M}
\newcommand{\bN}{\b N}
\newcommand{\bO}{\b O}
\newcommand{\bP}{\b P}
\newcommand{\bQ}{\b Q}
\newcommand{\bR}{\b R}
\newcommand{\bS}{\b S}
\newcommand{\bT}{\b T}
\newcommand{\bU}{\b U}
\newcommand{\bV}{\b V}
\newcommand{\bW}{\b W}
\newcommand{\bX}{\b X}
\newcommand{\bY}{\b Y}
\newcommand{\bZ}{\b Z}

\newcommand{\ba}{\b a}
\newcommand{\bb}{\b b}
\newcommand{\bc}{\b c}
\newcommand{\bd}{\b d}
\newcommand{\be}{\b e}
\newcommand{\bff}{\b f}
\newcommand{\bg}{\b g}
\newcommand{\bh}{\b h}
\newcommand{\bi}{\b i}
\newcommand{\bj}{\b j}
\newcommand{\bk}{\b k}
\newcommand{\bl}{\b l}
\newcommand{\bm}{\b m}
\newcommand{\bn}{\b n}
\newcommand{\bo}{\b o}
\newcommand{\bp}{\b p}
\newcommand{\bq}{\b q}
\newcommand{\br}{\b r}
\newcommand{\bs}{\b s}
\newcommand{\bt}{\b t}
\newcommand{\bu}{\b u}
\newcommand{\bv}{\b v}
\newcommand{\bw}{\b w}
\newcommand{\bx}{\b x}
\newcommand{\by}{\b y}
\newcommand{\bz}{\b z}

\newcommand{\cA}{\c A}
\newcommand{\cB}{\c B}
\newcommand{\cC}{\c C}
\newcommand{\cD}{\c D}
\newcommand{\cE}{\c E}
\newcommand{\cF}{\c F}
\newcommand{\cG}{\c G}
\newcommand{\cH}{\c H}
\newcommand{\cI}{\c I}
\newcommand{\cJ}{\c J}
\newcommand{\cK}{\c K}
\newcommand{\cL}{\c L}
\newcommand{\cM}{\c M}
\newcommand{\cN}{\c N}
\newcommand{\cO}{\c O}
\newcommand{\cP}{\c P}
\newcommand{\cQ}{\c Q}
\newcommand{\cR}{\c R}
\newcommand{\cS}{\c S}
\newcommand{\cT}{\c T}
\newcommand{\cU}{\c U}
\newcommand{\cV}{\c V}
\newcommand{\cW}{\c W}
\newcommand{\cX}{\c X}
\newcommand{\cY}{\c Y}
\newcommand{\cZ}{\c Z}

\newcommand{\dA}{\d A}
\newcommand{\dB}{\d B}
\newcommand{\dC}{\d C}
\newcommand{\dD}{\d D}
\newcommand{\dE}{\d E}
\newcommand{\dF}{\d F}
\newcommand{\dG}{\d G}
\newcommand{\dH}{\d H}
\newcommand{\dI}{\d I}
\newcommand{\dJ}{\d J}
\newcommand{\dK}{\d K}
\newcommand{\dL}{\d L}
\newcommand{\dM}{\d M}
\newcommand{\dN}{\d N}
\newcommand{\dO}{\d O}
\newcommand{\dP}{\d P}
\newcommand{\dQ}{\d Q}
\newcommand{\dR}{\d R}
\newcommand{\dS}{\d S}
\newcommand{\dT}{\d T}
\newcommand{\dU}{\d U}
\newcommand{\dV}{\d V}
\newcommand{\dW}{\d W}
\newcommand{\dX}{\d X}
\newcommand{\dY}{\d Y}
\newcommand{\dZ}{\d Z}
\newcommand{\fA}{\f A}
\newcommand{\fB}{\f B}
\newcommand{\fC}{\f C}
\newcommand{\fD}{\f D}
\newcommand{\fE}{\f E}
\newcommand{\fF}{\f F}
\newcommand{\fG}{\f G}
\newcommand{\fH}{\f H}
\newcommand{\fI}{\f I}
\newcommand{\fJ}{\f J}
\newcommand{\fK}{\f K}
\newcommand{\fL}{\f L}
\newcommand{\fM}{\f M}
\newcommand{\fN}{\f N}
\newcommand{\fO}{\f O}
\newcommand{\fP}{\f P}
\newcommand{\fQ}{\f Q}
\newcommand{\fR}{\f R}
\newcommand{\fS}{\f S}
\newcommand{\fT}{\f T}
\newcommand{\fU}{\f U}
\newcommand{\fV}{\f V}
\newcommand{\fW}{\f W}
\newcommand{\fX}{\f X}
\newcommand{\fY}{\f Y}
\newcommand{\fZ}{\f Z}

\newcommand{\fa}{\f a}
\newcommand{\fb}{\f b}
\newcommand{\fc}{\f c}
\newcommand{\fd}{\f d}
\newcommand{\fe}{\f e}
\newcommand{\ff}{\f f}
\newcommand{\fg}{\f g}
\newcommand{\fh}{\f h}
\newcommand{\fii}{\f i}
\newcommand{\fj}{\f j}
\newcommand{\fk}{\f k}
\newcommand{\fl}{\f l}
\newcommand{\fm}{\f m}
\newcommand{\fn}{\f n}
\newcommand{\fo}{\f o}
\newcommand{\fp}{\f p}
\newcommand{\fq}{\f q}
\newcommand{\fr}{\f r}
\newcommand{\fs}{\f s}
\newcommand{\ft}{\f t}
\newcommand{\fu}{\f u}
\newcommand{\fv}{\f v}
\newcommand{\fw}{\f w}
\newcommand{\fx}{\f x}
\newcommand{\fy}{\f y}
\newcommand{\fz}{\f z}

\newcommand{\rA}{\r A}
\newcommand{\rB}{\r B}
\newcommand{\rC}{\r C}
\newcommand{\rD}{\r D}
\newcommand{\rE}{\r E}
\newcommand{\rF}{\r F}
\newcommand{\rG}{\r G}
\newcommand{\rH}{\r H}
\newcommand{\rI}{\r I}
\newcommand{\rJ}{\r J}
\newcommand{\rK}{\r K}
\newcommand{\rL}{\r L}
\newcommand{\rM}{\r M}
\newcommand{\rN}{\r N}
\newcommand{\rO}{\r O}
\newcommand{\rP}{\r P}
\newcommand{\rQ}{\r Q}
\newcommand{\rR}{\r R}
\newcommand{\rS}{\r S}
\newcommand{\rT}{\r T}
\newcommand{\rU}{\r U}
\newcommand{\rV}{\r V}
\newcommand{\rW}{\r W}
\newcommand{\rX}{\r X}
\newcommand{\rY}{\r Y}
\newcommand{\rZ}{\r Z}

\newcommand{\ra}{\r a}
\newcommand{\rb}{\r b}
\newcommand{\rc}{\r c}
\newcommand{\rd}{\r d}
\newcommand{\re}{\r e}
\newcommand{\rf}{\r f}
\newcommand{\rg}{\r g}
\newcommand{\rh}{\r h}
\newcommand{\ri}{\r i}
\newcommand{\rj}{\r j}
\newcommand{\rk}{\r k}
\newcommand{\rl}{\r l}
\renewcommand{\rm}{\r m}
\newcommand{\rn}{\r n}
\newcommand{\ro}{\r o}
\newcommand{\rp}{\r p}
\renewcommand{\rq}{\r q}
\newcommand{\rr}{\r r}
\newcommand{\rs}{\r s}
\newcommand{\rt}{\r t}
\newcommand{\ru}{\r u}
\newcommand{\rv}{\r v}
\newcommand{\rw}{\r w}
\newcommand{\rx}{\r x}
\newcommand{\ry}{\r y}
\newcommand{\rz}{\r z}

\newcommand{\sA}{\s A}
\newcommand{\sB}{\s B}
\newcommand{\sC}{\s C}
\newcommand{\sD}{\s D}
\newcommand{\sE}{\s E}
\newcommand{\sF}{\s F}
\newcommand{\sG}{\s G}
\newcommand{\sH}{\s H}
\newcommand{\sI}{\s I}
\newcommand{\sJ}{\s J}
\newcommand{\sK}{\s K}
\newcommand{\sL}{\s L}
\newcommand{\sM}{\s M}
\newcommand{\sN}{\s N}
\newcommand{\sO}{\s O}
\newcommand{\sP}{\s P}
\newcommand{\sQ}{\s Q}
\newcommand{\sR}{\s R}
\newcommand{\sS}{\s S}
\newcommand{\sT}{\s T}
\newcommand{\sU}{\s U}
\newcommand{\sV}{\s V}
\newcommand{\sW}{\s W}
\newcommand{\sX}{\s X}
\newcommand{\sY}{\s Y}
\newcommand{\sZ}{\s Z}

\newcommand{\sfA}{\sf A}
\newcommand{\sfB}{\sf B}
\newcommand{\sfC}{\sf C}
\newcommand{\sfD}{\sf D}
\newcommand{\sfE}{\sf E}
\newcommand{\sfF}{\sf F}
\newcommand{\sfG}{\sf G}
\newcommand{\sfH}{\sf H}
\newcommand{\sfI}{\sf I}
\newcommand{\sfJ}{\sf J}
\newcommand{\sfK}{\sf K}
\newcommand{\sfL}{\sf L}
\newcommand{\sfM}{\sf M}
\newcommand{\sfN}{\sf N}
\newcommand{\sfO}{\sf O}
\newcommand{\sfP}{\sf P}
\newcommand{\sfQ}{\sf Q}
\newcommand{\sfR}{\sf R}
\newcommand{\sfS}{\sf S}
\newcommand{\sfT}{\sf T}
\newcommand{\sfU}{\sf U}
\newcommand{\sfV}{\sf V}
\newcommand{\sfW}{\sf W}
\newcommand{\sfX}{\sf X}
\newcommand{\sfY}{\sf Y}
\newcommand{\sfZ}{\sf Z}

\newcommand{\sfa}{\sf a}
\newcommand{\sfb}{\sf b}
\newcommand{\sfc}{\sf c}
\newcommand{\sfd}{\sf d}
\newcommand{\sfe}{\sf e}
\newcommand{\sff}{\sf f}
\newcommand{\sfg}{\sf g}
\newcommand{\sfh}{\sf h}
\newcommand{\sfi}{\sf i}
\newcommand{\sfj}{\sf j}
\newcommand{\sfk}{\sf k}
\newcommand{\sfl}{\sf l}
\newcommand{\sfm}{\sf m}
\newcommand{\sfn}{\sf n}
\newcommand{\sfo}{\sf o}
\newcommand{\sfp}{\sf p}
\newcommand{\sfq}{\sf q}
\newcommand{\sfr}{\sf r}
\newcommand{\sfs}{\sf s}
\newcommand{\sft}{\sf t}
\newcommand{\sfu}{\sf u}
\newcommand{\sfv}{\sf v}
\newcommand{\sfw}{\sf w}
\newcommand{\sfx}{\sf x}
\newcommand{\sfy}{\sf y}
\newcommand{\sfz}{\sf z}

\newcommand{\tA}{\mathtt{A}}
\newcommand{\tB}{\mathtt{B}}
\newcommand{\tC}{\mathtt{C}}
\newcommand{\tD}{\mathtt{D}}
\newcommand{\tE}{\mathtt{E}}
\newcommand{\tF}{\mathtt{F}}
\newcommand{\tG}{\mathtt{G}}
\newcommand{\tH}{\mathtt{H}}
\newcommand{\tI}{\mathtt{I}}
\newcommand{\tJ}{\mathtt{J}}
\newcommand{\tK}{\mathtt{K}}
\newcommand{\tL}{\mathtt{L}}
\newcommand{\tM}{\mathtt{M}}
\newcommand{\tN}{\mathtt{N}}
\newcommand{\tO}{\mathtt{O}}
\newcommand{\tP}{\mathtt{P}}
\newcommand{\tQ}{\mathtt{Q}}
\newcommand{\tR}{\mathtt{R}}
\newcommand{\tS}{\mathtt{S}}
\newcommand{\tT}{\mathtt{T}}
\newcommand{\tU}{\mathtt{U}}
\newcommand{\tV}{\mathtt{V}}
\newcommand{\tW}{\mathtt{W}}
\newcommand{\tX}{\mathtt{X}}
\newcommand{\tY}{\mathtt{Y}}
\newcommand{\tZ}{\mathtt{Z}}

\newcommand{\ta}{\mathtt{a}}
\newcommand{\tb}{\mathtt{b}}
\newcommand{\tc}{\mathtt{c}}
\newcommand{\td}{\mathtt{d}}
\newcommand{\te}{\mathtt{e}}
\newcommand{\tf}{\mathtt{f}}
\newcommand{\tg}{\mathtt{g}}
%\newcommand{\th}{\mathtt{h}}
\newcommand{\ti}{\mathtt{i}}
\newcommand{\tj}{\mathtt{j}}
\newcommand{\tk}{\mathtt{k}}
\newcommand{\tl}{\mathtt{l}}
\newcommand{\tm}{\mathtt{m}}
\newcommand{\tn}{\mathtt{n}}
%\newcommand{\to}{\mathtt{o}}
%\newcommand{\tp}{\mathtt{p}}
\newcommand{\tq}{\mathtt{q}}
%\newcommand{\tr}{\mathtt{r}}
\newcommand{\ts}{\mathtt{s}}
\renewcommand{\tt}{\mathtt{t}}
\newcommand{\tu}{\mathtt{u}}
\newcommand{\tv}{\mathtt{v}}
\newcommand{\tw}{\mathtt{w}}
\newcommand{\tx}{\mathtt{x}}
\newcommand{\ty}{\mathtt{y}}
\newcommand{\tz}{\mathtt{z}}

\newcommand{\CC}{\mathbf{C}}
\newcommand{\QQ}{\mathbf{Q}}
\newcommand{\ZZ}{\mathbf{Z}}
\newcommand{\RR}{\mathbf{R}}
%\newcommand{\FF}{\mathbf{F}}






%%%%%%%%%%%%%%%%%%%%%%%%%%%%%%%%%%%%%%%


\newcommand{\balpha}{\boldsymbol{\alpha}}
\newcommand{\biota}{\boldsymbol{\iota}}
\newcommand{\bkappa}{\boldsymbol{\kappa}}
\newcommand{\blambda}{\boldsymbol{\lambda}}
\newcommand{\bmu}{\boldsymbol{\mu}}
\newcommand{\btau}{\boldsymbol{\tau}}
\newcommand{\bbi}{\boldsymbol{i}}
\newcommand{\bbj}{\boldsymbol{j}}
\newcommand{\bbt}{\boldsymbol{t}}
\newcommand{\bbz}{\boldsymbol{z}}
\newcommand{\bbL}{\boldsymbol{L}}
\newcommand{\bbM}{\boldsymbol{M}}
\newcommand{\bbT}{\boldsymbol{T}}
\newcommand{\bbV}{\boldsymbol{V}}
\newcommand{\bbW}{\boldsymbol{W}}
\newcommand{\obj}{\text{\Flatsteel}}
\newcommand{\qbinom}[2]{\genfrac{[}{]}{0pt}{}{#1}{#2}}

\newcommand{\pres}[1]{{}^{#1}}
\newcommand{\floor}[1]{\lfloor{#1}\rfloor}
\newcommand{\ceil}[1]{\lceil{#1}\rceil}
\newcommand{\tube}[1]{\{\{{#1}\}\}}
\newcommand{\ab}{\r{ab}}
\newcommand{\abs}{\r{abs}}
%\newcommand{\ad}{\r{ad}}
\newcommand{\AH}{\r{AH}}
\newcommand{\bad}{\r{bad}}
\newcommand{\can}{\r{can}}
\newcommand{\CF}{\mathbbm{1}}
\newcommand{\cInd}{\rc\text{-}\r{Ind}}
\newcommand{\cl}{\r{cl}}
\newcommand{\Cl}{\r{Cl}}
\newcommand{\cont}{\r{cont}}
\newcommand{\cosp}{\r{cosp}}
\newcommand{\cris}{\r{cris}}
\newcommand{\cusp}{\r{cusp}}
\newcommand{\cyc}{\r{cyc}}
\newcommand{\dege}{\r{deg}}
\newcommand{\dgn}{\r{dgn}}
\newcommand{\df}{\r{def}}
\newcommand{\dr}{\r{dR}}
\newcommand{\et}{{\acute{\r{e}}\r{t}}}
\newcommand{\exc}{\r{exc}}
\newcommand{\ext}{\r{ext}}
\newcommand{\fin}{\r{fin}}
\newcommand{\FL}{\r{FL}}
\newcommand{\flg}{{\rf. \rl. }}
\newcommand{\free}{\r{free}}
\newcommand{\gr}{\r{gr}}
%\renewcommand{\graph}{\triangle}
\newcommand{\Gys}{\mathrm{Gys}}
\newcommand{\Hk}{\r{Hk}}
\newcommand{\HOM}{\c{H}om}
\newcommand{\HT}{\r{HT}}
\newcommand{\id}{\r{id}}
\newcommand{\inc}{\r{inc}}
\newcommand{\Inc}{\r{Inc}}
\newcommand{\inv}{\r{inv}}
\newcommand{\loc}{\r{loc}}
\newcommand{\lr}{\r{lr}}
\newcommand{\MF}{\sM\!\!\sF\!}
\newcommand{\mix}{\r{mix}}
\newcommand{\mnm}{\r{min}}
\newcommand{\mot}{\r{mot}}
\newcommand{\Nilp}{\r{Nil}}
\newcommand{\no}{\r{no}}
\newcommand{\ns}{\r{ns}}
\newcommand{\ordi}{\r{ord}}
\newcommand{\prim}{\r{prim}}
\newcommand{\ram}{\r{ram}}
\newcommand{\rat}{\r{rat}}
\newcommand{\RES}{\r{res}}
\newcommand{\Sat}{\r{Sat}}
\newcommand{\sh}{\r{sh}}
\newcommand{\sg}{\r{sg}}
\newcommand{\sing}{\r{sing}}
\renewcommand{\sp}{\r{sp}}
\newcommand{\sph}{\r{sph}}
\newcommand{\sep}{\r{sep}}
\newcommand{\ssp}{\r{ssp}}
\newcommand{\std}{\r{std}}
\newcommand{\St}{\r{St}}
\newcommand{\sym}{\r{sym}}
\newcommand{\tor}{\r{tor}}
\newcommand{\univ}{\r{univ}}
\newcommand{\unr}{\r{unr}}
\newcommand{\ur}{\r{ur}}
\newcommand{\pr}{\mathrm{pr}}
\newcommand{\even}{\mathrm{even}}
\newcommand{\eff}{\mathrm{eff}}
\newcommand{\odd}{\r{odd}}
\newcommand{\xra}{\xrightarrow}
\newcommand{\hra}{\hookrightarrow}

\newcommand{\Mod}{\sf{Mod}}
\newcommand{\Set}{\sf{Set}}
\newcommand{\Fun}{\sf{Fun}}
\newcommand{\Sch}{\sf{Sch}}
\newcommand{\defin}{\r{def}}
\newcommand{\indef}{\r{indef}}


\DeclareMathOperator{\ad}{ad}
\DeclareMathOperator{\AJ}{AJ}
\DeclareMathOperator{\Alb}{Alb}
\DeclareMathOperator{\As}{As}
\DeclareMathOperator{\Aut}{Aut}
\DeclareMathOperator{\AV}{AV}
\DeclareMathOperator{\BC}{BC}
\DeclareMathOperator{\CH}{CH}
\DeclareMathOperator{\Char}{char}
\DeclareMathOperator{\coker}{coker}
\DeclareMathOperator{\cond}{cond}
\DeclareMathOperator{\Corr}{Corr}
\DeclareMathOperator{\Def}{Def}
\DeclareMathOperator{\diag}{diag}
\DeclareMathOperator{\disc}{disc}
\DeclareMathOperator{\Disc}{Disc}
\DeclareMathOperator{\Div}{Div}
\DeclareMathOperator{\DIV}{div}
\DeclareMathOperator{\DL}{DL}
\DeclareMathOperator{\End}{End}
\DeclareMathOperator{\Ext}{Ext}
\DeclareMathOperator{\Fil}{Fil}
\DeclareMathOperator{\Fitt}{Fitt}
\DeclareMathOperator{\Fr}{Fr}
\DeclareMathOperator{\Frac}{Frac}
\DeclareMathOperator{\Frob}{Frob}
\DeclareMathOperator{\Gal}{Gal}
\DeclareMathOperator{\gap}{gap}
\DeclareMathOperator{\GL}{GL}
\DeclareMathOperator{\Gr}{Gr}
\DeclareMathOperator{\GU}{GU}
\DeclareMathOperator{\Hom}{Hom}
\DeclareMathOperator{\IP}{Im}
\DeclareMathOperator{\IM}{im}
\DeclareMathOperator{\ind}{ind}
\DeclareMathOperator{\Ind}{Ind}
\DeclareMathOperator{\IC}{IC}
\DeclareMathOperator{\IH}{IH}
\DeclareMathOperator{\Irr}{Irr}
\DeclareMathOperator{\Iso}{Iso}
\DeclareMathOperator{\Jac}{Jac}
\DeclareMathOperator{\Ker}{ker}
\DeclareMathOperator{\Lat}{Lat}
\DeclareMathOperator{\length}{length}
\DeclareMathOperator{\Lie}{Lie}
\DeclareMathOperator{\Map}{Map}
\DeclareMathOperator{\Mat}{Mat}
\DeclareMathOperator{\Max}{Max}
\DeclareMathOperator{\modulo}{mod}
\DeclareMathOperator{\Nm}{Nm}
\DeclareMathOperator{\Ob}{Ob}
\DeclareMathOperator{\ord}{ord}
\DeclareMathOperator{\Perv}{Perv}
\DeclareMathOperator{\PGL}{PGL}
\DeclareMathOperator{\Pic}{Pic}
\DeclareMathOperator{\rank}{rank}
\DeclareMathOperator{\RE}{Re}
%\DeclareMathOperator{\Res}{Res}
\DeclareMathOperator{\Sel}{Sel}
\DeclareMathOperator{\Stub}{Stub}
\DeclareMathOperator{\Sh}{Sh}
\DeclareMathOperator{\SL}{SL}
\DeclareMathOperator{\SO}{SO}
\DeclareMathOperator{\Sp}{Sp}
\DeclareMathOperator{\SP}{SP}
\DeclareMathOperator{\Span}{Span}
\DeclareMathOperator{\Spec}{Spec}
\DeclareMathOperator{\Spf}{Spf}
\DeclareMathOperator{\Std}{Std}
\DeclareMathOperator{\supp}{supp}
\DeclareMathOperator{\Sym}{Sym}
\DeclareMathOperator{\Tor}{Tor}
\DeclareMathOperator{\Tr}{Tr}
\DeclareMathOperator{\tr}{tr}
\DeclareMathOperator{\UG}{U}
\DeclareMathOperator{\val}{val}
\DeclareMathOperator{\Ver}{Ver}
%\DeclareMathOperator{\vol}{vol}
\DeclareMathOperator{\WD}{WD}
\DeclareMathOperator{\Art}{Art}


\def\var{\varphi}

\DeclareMathSymbol{\twoheadrightarrow} {\mathrel}{AMSa}{"10}
%\DeclareMathOperator*{\Hom}{Hom}
\DeclareMathOperator*{\esssup}{ess\, sup}
%\DeclareMathOperator*{\supp}{supp}
\DeclareMathOperator*{\dist}{dist}
%\DeclareMathOperator*{\loc}{loc}
\DeclareMathOperator*{\grad}{grad}
%\DeclareMathOperator*{\Div}{div}
%\DeclareMathOperator*{\tr}{tr}
%\DeclareMathOperator*{\id}{id}
%\DeclareMathOperator*{\rank}{rank}
%\DeclareMathOperator*{\Ker}{Ker}
\DeclareMathOperator*{\Diff}{Diff}
\DeclareMathOperator*{\sign}{sign}
\DeclareMathOperator*{\m}{mod}
\DeclareMathOperator*{\Scal}{Scal}
%\DeclareMathOperator*{\Ric}{Ric}
\DeclareMathOperator*{\Sec}{Sec}
\DeclareMathOperator*{\Hess}{Hess}
%math
\newcommand{\norm}[1]{\left\lVert#1\right\rVert}%norm%
%\newcommand{\abs}[1]{\left\lvert#1\right\rvert}%absolute%
\newcommand{\set}[1]{\left\{#1\right\}}%set%
\newcommand{\hin}[2]{\left\langle#1, #2\right\rangle}%Hermite inner product%
\newcommand{\rin}[2]{\left(#1, #2\right)}%the real part of Hermite inner product%

\newcommand{\field}[1]{\mathbb{#1}}%field%
%\newcommand{\R}{\field{R}}
\newcommand{\Com}{\field{C}}
\newcommand{\Quat}{\field{H}}
\newcommand{\Lg}[1]{\mathrm{#1}}%Lie group%
%\newcommand{\Cl}{\mathbb{C}l}
\newcommand{\La}[1]{\mathfrak{#1}}%Lie algebra%
\newcommand{\tensor}[1]{\mathsf{#1}}%tensor%
\newcommand{\vect}[1]{\mathbf{#1}}%vector%
\newcommand{\To}{\longrightarrow}
\newcommand{\rmn}[1]{\romannumeral#1}%roman number%
\newcommand{\Rmn}[1]{\uppercase\expandafter{\romannueral#1}}%uppercase roman number%


\def\RR{\mathbb{R}}
\def\CC{\mathbb{C}}
\def\ii{\sqrt{-1}}
\def\NN{\mathcal{N}}
%\def\cc{\widetilde{\nabla }}
\def\ttau{\widetilde{\tau }}
\def\grad{\mathrm{grad\, }}
\def\Hess{\mathrm{Hess\, }}
\def\Ric{\mathrm{Ric\, }}
\def\Div{\mathrm{div\, }}
\def\Ker{\mathrm{Ker\, }}
\def\Im{\mathrm{Im\, }}
\def\tr{\mathrm{tr\, }}
%\renewcommand{\thethm}{\Alph{thm}}
\numberwithin{equation}{subsection}


\makeatletter
\newcommand{\fakephantomsection}{%
	\Hy@MakeCurrentHref{\@currenvir. \the\Hy@linkcounter}
	\Hy@raisedlink{\hyper@anchorstart{\@currentHref}\hyper@anchorend}%
}
\makeatother

\def\calA{\mathcal{A}}
\def\calB{\mathcal{B}}
\def\calC{\mathcal{C}}
\def\calD{\mathcal{D}}
\def\calE{\mathcal{E}}
\def\calF{\mathcal{F}}
\def\calG{\mathcal{G}}
\def\calH{\mathcal{H}}
\def\calL{\mathcal{L}}
\def\calM{\mathcal{M}}
\def\calN{\mathcal{N}}
\def\calO{\mathcal{O}}
\def\calR{\mathcal{R}}
\def\calS{\mathcal{S}}
\def\calT{\mathcal{T}}
\def\calU{\mathcal{U}}
\def\calV{\mathcal{V}}
\def\calW{\mathcal{W}}
\def\calX{\mathcal{X}}
\def\calY{\mathcal{Y}}
\def\calZ{\mathcal{Z}}
\def\gothA{\mathfrak{A}}
\def\gothC{\mathfrak{C}}
\def\gothD{\mathfrak{D}}
\def\gothE{\mathfrak{E}}
\def\gothm{\mathfrak{m}}
\def\gotho{\mathfrak{o}}
\def\gothp{\mathfrak{p}}
\def\gothM{\mathfrak{M}}
\def\gothN{\mathfrak{N}}
\def\gothP{\mathfrak{P}}
\def\gothQ{\mathfrak{Q}}
\def\gothR{\mathfrak{R}}
\def\gothS{\mathfrak{S}}
\def\gothH{\mathfrak{H}}
\def\gothT{\mathfrak{T}}
\def\gothU{\mathfrak{U}}
\def\gothV{\mathfrak{V}}
\def\gothX{\mathfrak{X}}
\def\gothZ{\mathfrak{Z}}
\def\bfA{\mathbf{A}}
\def\bfB{\mathbf{B}}
\def\bfI{\mathbf{I}}
\def\AAA{\mathbb{A}}
\def\BB{\mathbb{B}}
\def\CC{\mathbb{C}}
\def\DD{\mathbb{D}}
\def\FF{\mathbb{F}}
\def\GG{\mathbb{G}}
\def\HH{\mathbb{H}}
\def\II{\mathbb{I}}
\def\NN{\mathbb{N}}
\def\PP{\mathbb{P}}
\def\QQ{\mathbb{Q}}
\def\RR{\mathbb{R}}
\def\TT{\mathbb{T}}
\def\XX{\mathbb{X}}
\def\ZZ{\mathbb{Z}}
\def\bfa{\mathbf{a}}
\def\bfb{\mathbf{b}}
\def\bfe{\mathbf{e}}
\def\bff{\mathbf{f}}
\def\bfg{\mathbf{g}}
\def\bfk{\mathbf{k}}
\def\bfm{\mathbf{m}}
\def\bfv{\mathbf{v}}
\def\bfw{\mathbf{w}}
\def\bfx{\mathbf{x}}
\def\bfD{\mathbf{D}}
\def\bfL{\mathbf{L}}
\def\bfT{\mathbf{T}}
\def\bfC{\mathbf{C}}
\def\bfE{\mathbf{E}}
\def\bfM{\mathbf{M}}
\def\bfV{\mathbf{V}}
\def\bfP{\mathbf{P}}
\def\bfR{\mathbf{R}}
\def\rmm{\mathrm{m}}
\def\rmn{\mathrm{n}}
\def\rmt{\mathrm{t}}
\def\rmA{\mathrm{A}}
\def\rmB{\mathrm{B}}
\def\rmC{\mathrm{C}}
\def\rmD{\mathrm{D}}
\def\rmG{\mathrm{G}}
\def\rmH{\mathrm{H}}
\def\rmI{\mathrm{I}}
\def\rmK{\mathrm{K}}
\def\rmL{\mathrm{L}}
\def\rmM{\mathrm{M}}
\def\rmP{\mathrm{P}}
\def\rmU{\mathrm{U}}
\def\rmR{\mathrm{R}}
\def\rmT{\mathrm{T}}
\def\rmY{\mathrm{Y}}
\def\rmZ{\mathrm{Z}}
\def\scrB{\mathscr{B}}
\def\scrC{\mathscr{C}}
\def\i{\mathbf{i}}
\def\j{\mathbf{j}}
\def\k{\mathbf{k}}
\def\rmS{\mathrm{S}}
\def\sfS{\mathsf{S}}
\def\sfc{\mathsf{c}}
\def\sfw{\mathsf{w}}
\def\sfz{\mathsf{z}}
\def\tti{\mathtt{i}}
\def\frakf{\mathfrak{f}}

%\newcommand{\et}{\mathrm{\acute et}}
%\newcommand{\Proj}{\mathrm{Proj}}
%\newcommand{\proj}{\mathrm{proj}}
%\DeclareMathOperator{\Cone}{Cone}
%\DeclareMathOperator{\Coker}{Coker}
%\DeclareMathOperator{\Def}{Def}
%\DeclareMathOperator{\End}{End}
%\DeclareMathOperator{\Ext}{Ext}
%\DeclareMathOperator{\Gal}{Gal}
%\DeclareMathOperator{\GL}{GL}
%%\DeclareMathOperator{\Ker}{Ker}
%%\DeclareMathOperator{\Hom}{Hom}
%\DeclareMathOperator{\HT}{HT}
%\DeclareMathOperator{\HP}{HP}
%\DeclareMathOperator{\Lie}{Lie}
%\DeclareMathOperator{\rank}{rank}
%\DeclareMathOperator{\Max}{Max}
%\DeclareMathOperator{\NCL}{NCL}
%\DeclareMathOperator{\NP}{NP}
%\DeclareMathOperator{\Res}{Res}
%\DeclareMathOperator{\Spec}{Spec}
%\DeclareMathOperator{\Spf}{Spf}
\DeclareMathOperator{\Spm}{Spm}
%\DeclareMathOperator{\Sym}{Sym}
%\DeclareMathOperator{\Tor}{Tor}
%\newcommand{\op}{\mathrm{op}}
%\newcommand{\Dig}{\mathrm{Dig}}
%\newcommand{\Sat}{\mathrm{Sat}}
%\newcommand{\unit}{\mathrm{unit}}
%\newcommand{\AL}{\mathrm{AL}}
%\newcommand{\an}{\mathrm{an}}
%\newcommand{\Berk}{\mathrm{Berk}}
%\newcommand{\bbsigma}{\boldsymbol{\sigma}}
%\newcommand{\can}{\mathrm{can}}
%\newcommand{\cont}{\mathrm{cont}}
%\newcommand{\cInd}{\textrm{c-Ind}}
\newcommand{\cycl}{\mathrm{cycl}}
\newcommand{\crys}{\mathrm{crys}}
\newcommand{\pcrys}{\mathrm{pcrys}}
%\newcommand{\dR}{\mathrm{dR}}
%\newcommand{\dual}{\vee}
%\newcommand{\Fil}{\mathrm{Fil}}
%\newcommand{\fin}{\mathrm{fin}}
%\newcommand{\Fr}{\mathrm{Fr}}
%\newcommand{\Hol}{\mathrm{Hol}}
%\newcommand{\HOL}{\mathrm{HOL}}
%%\newcommand{\id}{\mathrm{id}}
%\newcommand{\Mod}{\mathrm{Mod}}
\newcommand{\new}{\mathrm{new}}
\newcommand{\old}{\textrm{-}\mathrm{old}}
\newcommand{\sm}{\mathrm{sm}}
\newcommand{\Iw}{\mathrm{Iw}}
\renewcommand{\det}{\mathrm{det}}
%\newcommand{\unr}{\mathrm{unr}}
%\newcommand{\nord}{\mathrm{nord}}
%\newcommand{\ord}{\mathrm{ord}}
\newcommand{\sgn}{\mathrm{sgn}}
%\newcommand{\Nm}{\mathrm{Nm}}
%\newcommand{\pr}{\mathrm{pr}}
%\DeclareMathOperator{\corank}{corank}
%\newcommand{\Qp}{\QQ_p}
%\newcommand{\Rep}{\mathrm{Rep}}
%\newcommand{\swap}{\mathrm{swap}}
%\newcommand{\slp}{\mathrm{slp}}
%\newcommand{\rig}{\mathrm{rig}}
%\newcommand{\Tate}{\mathrm{Tate}}
%\newcommand{\tame}{\mathrm{tame}}
%\newcommand{\tw}{\mathrm{tw}}
%\newcommand{\wt}{\mathrm{wt}}
%\newcommand{\WD}{\mathrm{WD}}
\newcommand{\WT}{\mathrm{WT}}
\newcommand{\Zp}{\ZZ_p}
%\newcommand{\SL}{\mathrm{SL}}
%\newcommand{\cl}{\mathrm{cl}}
\newcommand{\inte}{\mathrm{int}}
\newcommand{\Err}{\mathrm{Err}}
\newcommand{\lan}{\mathrm{l. an}}
\newcommand{\alg}{\mathrm{alg}}
\newcommand{\naive}{\mathrm{naive}}
\newcommand{\LP}{\mathrm{LP}}
%\newcommand{\Sh}{\mathrm{Sh}}
%\newcommand{\ur}{\mathrm{ur}}
%\newcommand{\Frob}{\mathrm{Frob}}
%\newcommand{\univ}{\mathrm{univ}}
\newcommand{\JH}{\mathrm{JH}}
\newcommand{\Mult}{\mathrm{Mult}}
\newcommand{\power}{\mathrm{power}}
\newcommand{\tri}{\mathrm{tri}}
\newcommand{\pro}{\mathrm{pro}}
\newcommand{\nS}{\mathrm{nS}}
\newcommand{\Vtx}{\mathrm{Vtx}}
\newcommand{\lead}{\mathrm{lead}}
\newcommand{\midd}{\mathrm{mid}}
%\newcommand{\even}{\mathrm{even}}
%\newcommand{\odd}{\mathrm{odd}}
\newcommand{\defect}{\mathrm{defect}}
\newcommand{\Mah}{\mathrm{Mah}}
\newcommand{\reg}{\mathrm{reg}}
\newcommand{\rec}{\mathrm{rec}}
\newcommand{\soc}{\mathrm{soc}}

%\DeclareMathOperator{\length}{length}
%\DeclareMathOperator{\Char}{Char}
\DeclareMathOperator{\Spa}{Spa}
\DeclareMathOperator{\Spc}{Spc}
\DeclareMathOperator{\Diag}{Diag}
\DeclareMathOperator{\Eig}{Eig}
%\DeclareMathOperator{\Ind}{Ind}
\allowdisplaybreaks




\def\pdz#1{\dfrac{\partial}{\partial z_{#1}}}
\def\pd#1{\dfrac{\partial}{\partial #1}}
\def\f#1#2{\frac{#1}{#2}}
\def\p#1#2{\dfrac{\partial #1}{\partial#2}}
\def\pa{\partial}
\def\dt{\frac{d}{dt}}
\def\na{\nabla}
\def\lap{\bigtriangleup}
\def\a{\alpha}
\def\b{\beta}
\def\ga{\gamma}
\def\vph{\varphi}
\def\({\left (}
\def\){\right )}
\def\<{\langle}
\def\>{\rangle}
\def\ma{|\mathring{A}|^2}
\def\ch{\mathring{h}_{ij}}
%\def\bint{-\!\!\!\!\!\!\!\int}

\newcommand{\bel}[1]{\begin{equation}\label{#1}}
	\newcommand{\lab}[1]{\label{#1}}
	\newcommand{\beq}{\begin{equation}}
		
		%\newcommand{\ba}{\begin{eqnarray}}
			\newcommand{\ea}{\end{eqnarray}}
		%	\newcommand{\rf}[1]{(\ref{#1})}
		\newcommand{\qe}{\end{equation}}
	\newcommand{\eeq}{\end{equation}}
\iffalse
\newcommand{\Yub}{\mathcal{Y}_{\ubar}}
\newcommand{\Lbar}{\underline{L}}
\newcommand{\Hbar}{\underline{H}}
\newcommand{\psig}{\psi_g}
\newcommand{\aln}{(\al \nabla)^i}
\newcommand{\ali}{a^{\frac{i}{2}}}
\newcommand{\scaletwoHprime}[1]{\lVert #1 \rVert_{L^2_{(sc)}(H_{u^{\prime}}^{(0, \ubar)}) }}
\newcommand{\twoSuprime}[1]{\lVert #1 \rVert_{L^2(S_{u^{\prime}, \ubar})}}
\newcommand{\scaletwoHbarprime}[1]{\lVert #1 \rVert_{L^2_{(sc)}(\Hbar_{\ubar^{\prime}}^{(u_{\infty}, u)}) }}
\newcommand{\tbeta}{\beta^R}
\newcommand{\nablaF}{\nabla^{i_5}}
\newcommand{\hnablaF}{\hnabla^{i_5}}
\newcommand{\tbetabar}{\betabar^R}
\newcommand{\hsp}{\hspace{. 5mm}}
\newcommand{\omegabar}{\underline{\omega}}
\newcommand{\gslash}{\slashed{g}}
\newcommand{\psibar}{\underline{\psi}}
\newcommand{\psiKl}{\psi_{\mathfrak{Kl}}}
\newcommand{\psibarKl}{\psibar_{\mathfrak{Kl}}}
\newcommand{\dubarprime}{\hspace{. 5mm} \text{d} \ubar^{\prime}}
\newcommand{\hnabla}{\nabla}
\newcommand{\intu}{\int_{u_{\infty}}^u}
\newcommand{\intubar}{\int_0^{\ubar}}
\newcommand{\del}{\partial}
\newcommand{\F}{{{F}^P}_Q}
\newcommand{\chihat}{\hat{\chi}}
\newcommand{\g}{\gamma}
\newcommand{\upr}{\lvert u^\prime \rvert}
\newcommand{\al}{a^{\frac{1}{2}}}
\newcommand{\chibar}{\underline{\chi}}
\newcommand{\chibarhat}{\underline{\hat{\chi}}}
\newcommand{\ubar}{\underline{u}}
\newcommand{\e}{\mathrm{e}}
\newcommand{\sumit}{\sum_{i_1+i_2+i_3=i}}
\newcommand{\sumitm}{\sum_{i_1+i_2+i_3=i-1}}
\newcommand{\sumiF}{\sum_{i_1+i_2+i_3+i_4+i_5=i}}
\newcommand{\sumiFi}{\sum_{i_1+i_2+i_3+i_4+i_5=i-1}}
\newcommand{\sumif}{\sum_{i_1+i_2+i_3+i_4=i}}
\newcommand{\sumifi}{\sum_{i_1+i_2+i_3+i_4=i-1}}
\newcommand{\sumifim}{\sum_{i_1+i_2+i_3+i_4=i-2}}
\newcommand{\sumjf}{\sum_{J_1+J_2+J_3+J_4=J}}
\newcommand{\sumjfj}{\sum_{J_1+J_2+J_3+J_4=J-1}}
\newcommand{\sumjfjm}{\sum_{J_1+J_2+J_3+J_4=J-2}}
%\newcommand{\tr}{\text{tr}}
\newcommand{\be}{\begin{equation}}
	\newcommand{\ee}{\end{equation}}
\newcommand{\dtext}{\text{d}}
\newcommand{\bm}{\begin{align}*}
	\newcommand{\enm}{\end{align}*}
\newcommand{\hess}{\operatorname{Hess}}
\newcommand{\supp}{\operatorname{supp}}
\newcommand{\scalep}[1]{\lVert{#1} \rVert_{L^p_{(sc)}(u, \underline{u})}}
\newcommand{\scaleinfinityS}[1]{\lVert{#1} \rVert_{\mathcal{L}^\infty_{(sc)}(S)}}
\newcommand{\scaleinfinitySu}[1]{\lVert{#1} \rVert_{L^\infty_{(sc)}(S_{u, \ubar})}}
\newcommand{\scaleinfinitySuprime}[1]{\lVert{#1} \rVert_{L^\infty_{(sc)}(S_{u^{\prime}, \ubar})}}
\newcommand{\ScaleinfinitySuprime}[1]{\big\lVert{#1} \big\rVert_{L^\infty_{(sc)}(S_{u^{\prime}, \ubar})}}
\newcommand{\scaleinfinitySuprimeubarprime}[1]{\lVert{#1} \rVert_{L^\infty_{(sc)}(S_{u^{\prime}, \ubar^\prime})}}
\newcommand{\ScaleinfinitySuprimeubarprime}[1]{\big\lVert{#1} \big\rVert_{L^\infty_{(sc)}(S_{u^{\prime}, \ubar^\prime})}}
\newcommand{\modulo}[1]{{\left|#1\right|}}

\newcommand{\Dist}{\operatorname{dist}}
\newcommand{\const}{\operatorname{const}}
\newcommand{\snr}[1]{\lvert #1\rvert}
\newcommand{\dv}{\operatorname{div}}
\newcommand{\vj}{v^{j}}
\newcommand{\tildetr}{\widetilde{\tr \chibar}}
\newcommand{\bx}{B_{r}(x_{0})}
\newcommand{\da}{D_{\alpha}}
\newcommand{\nr}[1]{\lVert #1 \rVert}
\newcommand{\gr}[1]{\mathbf{#1}}
\newcommand{\Lc}{\mathcal{L}}
\newcommand{\Y}{\mathrm{\Upsilon}}
\newcommand{\ubaruS}[3]{{L_{\ubar}^{#1} L_u^{#2} L_{\vphantom{\ubar}}^{#3}(S)}}
\newcommand{\nablap}{\nabla^{i_1}\psi_g^{i_2}}
\newcommand{\nablapp}{\nabla^{i_1}\psi_g^{i_2+1}}
\newcommand{\uubarS}[3]{{L_{u}^{#1} L_{\ubar}^{#2} L_{\vphantom{\ubar}}^{#3}(S_{u, \ubar})}}
\newcommand{\uubarSu}[3]{{L_{u}^{#1} L_{\ubar}^{#2} L_{\vphantom{\ubar}}^{#3}(S_{u, \ubar})}}
\newcommand{\scalefourS}[1]{\lVert{#1} \rVert_{L^4_{(sc)}(S)}}
\newcommand{\scaletwoS}[1]{\lVert{#1} \rVert_{L^2_{(sc)}(S)}}
\newcommand{\twoSu}[1]{\lVert{#1} \rVert_{L^2(S_{u, \ubar})}}
\newcommand{\pSu}[1]{\lVert{#1} \rVert_{L^p(S_{u, \ubar})}}
\newcommand{\inftySu}[1]{\lVert{#1} \rVert_{L^{\infty}(S_{u, \ubar})}}
\newcommand{\oneSu}[1]{\lVert{#1}\rVert_{L^1(S_{u, \ubar})}}
\newcommand{\fourSu}[1]{\lVert{#1} \rVert_{L^4(S_{u, \ubar})}}
\newcommand{\bespeq}{\begin{equation}\begin{split}}
		\newcommand{\espeq}{\end{split}\end{equation}}
\newcommand{\scaletwoSu}[1]{\lVert{#1} \rVert_{L^2_{(sc)}(S_{u, \ubar})}}
\newcommand{\scaleoneSu}[1]{\lVert{#1} \rVert_{L^1_{(sc)}(S_{u, \ubar})}}
\newcommand{\ScaletwoSu}[1]{\big\lVert{#1} \big\rVert_{L^2_{(sc)}(S_{u, \ubar})}}
\newcommand{\scaletwoSuprime}[1]{\lVert{#1} \rVert_{L^2_{(sc)}(S_{u^\prime, \ubar})}}
\newcommand{\scaletwoSuzprime}[1]{\lVert{#1} \rVert_{L^2_{(sc)}(S_{u^\prime, 0})}}
\newcommand{\ScaletwoSuprime}[1]{\big\lVert{#1} \big\rVert_{L^2_{(sc)}(S_{u^\prime, \ubar})}}
\newcommand{\scaleoneSuprimeubarprime}[1]{\lVert{#1} \rVert_{L^1_{(sc)}(S_{u^\prime, \ubar^\prime})}}
\newcommand{\scaletwoSuprimeubarprime}[1]{\lVert{#1} \rVert_{L^2_{(sc)}(S_{u^\prime, \ubar^\prime})}}
\newcommand{\ScaletwoSuprimeubarprime}[1]{\big\lVert{#1} \big\rVert_{L^2_{(sc)}(S_{u^\prime, \ubar^\prime})}}
\newcommand{\scaletwoSuubarprime}[1]{\lVert{#1} \rVert_{L^2_{(sc)}(S_{u, \ubar^\prime})}}
\newcommand{\scaletwoSuzubarprime}[1]{\lVert{#1} \rVert_{L^2_{(sc)}(S_{u_{\infty}, \ubar^\prime})}}
\newcommand{\ScaletwoSuubarprime}[1]{\Big \lVert{#1} \Big \rVert_{L^2_{(sc)}(S_{u, \ubar^\prime})}}
\newcommand{\LpSu}[1]{\lVert #1 \rVert_{L^p(\Suu)}}
\newcommand{\LinftySu}[1]{\lVert #1 \rVert_{L^{\infty}(\Suu)}}
\newcommand{\LpHu}[1]{\lVert #1 \rVert_{L^p(\Hu)}}
\newcommand{\LpHbaru}[1]{\lVert #1 \rVert_{L^p(\Hbu)}}
\newcommand{\scalefourH}[1]{\lVert{#1} \rVert_{\mathcal{L}^4_{(sc)}(H)}}
\newcommand{\alphabar}{\underline{\alpha}}
\newcommand{\betabar}{\underline{\beta}}
\newcommand{\etabar}{\underline{\eta}}
\newcommand{\scalefourHbar}[1]{\lVert{#1} \rVert_{L^4_{(sc)}(\underline{H})}}
\newcommand{\scaletwoH}[1]{\lVert{#1} \rVert_{L^2_{(sc)}(H)}}
\newcommand{\scaletwoHbar}[1]{\lVert{#1} \rVert_{L^2_{(sc)}(\underline{H})}}
\newcommand{\scalefourHu}[1]{\lVert{#1} \rVert_{L^4_{(sc)}(H_u^{(0, \underline{u})})}}
\newcommand{\scalefourHbaru}[1]{\lVert{#1} \rVert_{L^4_{(sc)}(\underline{H}_{\underline{u}}^{(0, u)})}}
\newcommand{\scaletwoHzero}[1]{\lVert{#1} \rVert_{L^2_{(sc)}(H_{u_{\infty}}^{(0, \underline{u})})}}
\newcommand{\scaletwoHu}[1]{\lVert{#1} \rVert_{L^2_{(sc)}(H_u^{(0, \underline{u})})}}
\newcommand{\scaletwoHbarzero}[1]{\lVert{#1} \rVert_{L^2_{(sc)}(\underline{H}_{0}^{(u_{\infty}, u)})}}
\newcommand{\scaletwoHbaru}[1]{\lVert{#1} \rVert_{L^2_{(sc)}(\underline{H}_{\underline{u}}^{(u_{\infty}, u)})}}
\newcommand{\dl}[1]{\epsilon #1}
\newcommand{\dt}[1]{\epsilon \tilde{#1}}
\newcommand{\scalepinitialHbar}[1]{\lVert{#1} \rVert_{L^p_{(sc)}(u, 0)}}
\newcommand{\scalepinitialH}[1]{\lVert{#1} \rVert_{L^p_{(sc)}(0, \underline{u})}}
%\newcommand{\br}{B_{r}(a)}
%\newcommand{\dbr}{\partial B_{r}(a)}
\newcommand{\Lloc}[1]{\mathbf{L^{#1}_{loc}}}
\newcommand{\C}[1]{\mathbf{C^{#1}}}
\newcommand{\Cc}[1]{\mathbf{C_c^{#1}}}
\newcommand{\Hutwo}[1]{\lVert #1 \rVert_{L^2(H_u)}}
\newcommand{\Hbarubartwo}[1]{\lVert #1 \rVert_{L^2(\Hbar_{\ubar})}}
\newcommand{\diver}{\operatorname{div}}
\newcommand{\omegad}{\omega^{\dagger}}
\newcommand{\omegabard}{\omegabar^{\dagger}}
\newcommand{\duprime}{\hspace{. 5mm} \text{d}u^{\prime}}
\newcommand{\mubar}{\underline{\mu}}
\newcommand{\kappabar}{\underline{\kappa}}
\newcommand{\sh}{\tau_{hs}}
\newcommand{\Suu}{S_{u, \ubar}}
\newcommand{\Suo}{S_{u, 0}}
\newcommand{\Sou}{S_{0, \ubar}}
\newcommand{\hs}{\tau_{-hs}}
\newcommand{\Dt}{D_{t}}
\newcommand{\ep}{\sh^{+}}
\newcommand{\eme}{\hs^{-}}
\newcommand{\utne}{\tilde{u}^{n}_{\varepsilon}}
\newcommand{\une}{u^{n}_{\varepsilon}}
\newcommand{\un}{u^{n}}
\newcommand{\ute}{\tilde{u}_{\varepsilon}}
\newcommand{\wu}{W^{1, p}_{u}}
\newcommand{\br}{B_{r}(0)}
\newcommand{\brp}{B_{r'}(0)}
\newcommand{\brs}{B_{r''}(0)}
\newcommand{\bs}{B_{s}(0)}
\newcommand{\bt}{B_{t}(0)}
\newcommand{\aupr}{\frac{a}{\upr^2}}
\newcommand{\haln}{(\al \hnabla)^i}
\newcommand{\RN}{\mathbb{R}^{N}}
\newcommand{\SN}{\mathbb{S}^{N-1}}
\newcommand{\N}{\mathbb{N}}
\newcommand{\gn}{G^{n}}
\newcommand{\nablat}{\nabla^{i_3}}
\newcommand{\hnablat}{\hnabla^{i_3}}
\newcommand{\nablaf}{\nabla^{i_4}}
\newcommand{\hnablaf}{\hnabla^{i_4}}
\newcommand{\Hodge}[1]{\prescript{*}{}{#1}}
\newcommand{\gne}{G^{n}_{\varepsilon}}
\newcommand{\ff}[1]{\int_{\br}F(D #1) \ dx}
\newcommand{\ffp}[1]{\int_{\brp}F(D #1) \ dx}
\newcommand{\intscaletwoSuubarprime}[1]{\int_0^{\ubar} \scaletwoSuubarprime{#1} \hspace{. 5mm} \text{d} \ubar^\prime}
\newcommand{\intscaleTwoSuubarprime}[1]{\int_0^{\ubar} \ScaletwoSuubarprime{#1} \hspace{. 5mm} \text{d} \ubar^\prime}
\newcommand{\fn}[1]{\int_{\br}F(D #1)+\frac{n^{2}}{4}\left(\snr{#1}^{2}-1\right)^{2}+\frac{1}{p}\snr{#1 -u}^{p}  \ dx}
\newcommand{\fne}[1]{\int_{\brp}F(D #1)+\frac{n^{2}}{4}\left(\snr{#1}^{2}-1\right)^{2}+\frac{1}{p}\snr{#1 - \ute}^{p}+\frac{\tilde{\varepsilon}}{2}\snr{D^{j} #1}^{2}  \ dx}
\newcommand{\fnp}[1]{\int_{\brp}F(D #1)+\frac{n^{2}}{4}\left(\snr{#1}^{2}-1\right)^{2}+\frac{1}{p}\snr{#1 - \ute}^{p}  \ dx}
\newcommand{\ffnp}[1]{\int_{\brp}F(D #1)+\frac{n^{2}}{4}\left(\snr{#1}^{2}-1\right)^{2}+\frac{1}{p}\snr{#1 - u}^{p}  \ dx}
\newcommand{\brl}{B_{\tilde{r}}(0)}
\newcommand{\gu}[1]{\left(n^{2}\left(\snr{#1}^{2}-1\right)#1 + \snr{#1 - \ute}^{p-2}(#1 - \ute) \right)}
\newcommand{\X}{\xi}
\newcommand{\E}{\eta}
\newcommand{\vn}{v_{j}}
\newcommand{\alphabarF}{\alphabar^F}
\newcommand{\rhoF}{\rho^F}
\newcommand{\sigmaF}{\sigma^F}
\newcommand{\alphaF}{\alpha^F}
\newcommand{\YM}{\mathcal{Y}}
\newcommand{\ub}{\underline{u}}
\newcommand{\Cb}{\underline{C}}
\newcommand{\Lb}{\underline{L}}
\newcommand{\Lh}{\hat{L}}
\newcommand{\Lbh}{\hat{\Lb}}
\newcommand{\phib}{\underline{\phi}}
\newcommand{\Phib}{\underline{\Phi}}
\newcommand{\Db}{\underline{D}}
\newcommand{\Dh}{\widehat{D}}
\newcommand{\Dbh}{\hat{\Db}}
\newcommand{\omb}{\underline{\omega}}
\newcommand{\omh}{\hat{\omega}}
\newcommand{\ombh}{\hat{\omb}}
\newcommand{\Pb}{\underline{P}}
\newcommand{\chib}{\underline{\chi}}
\newcommand{\chih}{\hat{\chi}}
\newcommand{\chibh}{\hat{\chib}}




\newcommand{\alb}{\underline{\alpha}}
\newcommand{\zeb}{\underline{\zeta}}
\newcommand{\beb}{\underline{\beta}}
\newcommand{\etb}{\underline{\eta}}
\newcommand{\Mb}{\underline{M}}
\newcommand{\oth}{\hat{\otimes}}




\def\a {\alpha}
\def\b {\beta}
\def\ab {\alphab}
\def\bb {\betab}
\def\nab {\nabla}





\def\ub {\underline{u}}
\def\th {\theta}
\def\Lb {\underline{L}}
\def\Hb {\underline{H}}
\def\chib {\underline{\chi}}
\def\chih {\hat{\chi}}
\def\chibh {\hat{\underline{\chi}}}
\def\omegab {\underline{\omega}}
\def\etab {\underline{\eta}}
\def\betab {\underline{\beta}}
\def\alphab {\underline{\alpha}}
\def\Psib {\underline{\Psi}}
\def\hot{\widehat{\otimes}}
\def\Phib {\underline{\Phi}}
\def\thb {\underline{\theta}}
\def\t {\tilde}
\def\st {\tilde{s}}







%%%%%%%%
%%%%%%%%
%\def\d {\epsilon}
\def\f {\frac}
\def\i {\infty}
\def\l {\bigg(}
\def\r {\bigg)}
\def\S {S_{u, \underline{u}}}
\def\o{\omega}
\def\O{\Omega}
%\def\be{\begin{equation}\begin{split}}
		%\def\en{\end{split}\end{equation}}

\def\od{\omega^{\dagger}}
\def\ombd{\underline{\omega}^{\dagger}}
\def\K{K-\frac{1}{|u|^2}}
\def\ut{\frac{1}{|u|^2}}
\def\Kb{K-\frac{1}{(u+\underline{u})^2}}
\def\M{\mathcal}
\def\p{\psi}

\def\D{\epsilon}
\def\T{\Theta}
\def\s{S_{u', \underline{u}'}}
\def\Hu{H_u^{(0, \underline{u})}}
\def\Hbu{\underline{H}_{\underline{u}}^{(u_{\infty}, u)}}
%\def\ee{(\eta, \underline{\eta})}

\def\at{a^{\f12}}
\def\sigmac{\check{\sigma}}
\def\p{\psi}
\def\q{\underline{\psi}}
\def\ls{\leq}
\def\de{\epsilon}
\def\ls{\lesssim}
\def\oo{\Omega\mbox{tr}\chib-\frac{2}{u}}
\def\om{\omega}
\def\Om{\Omega}
%%%%%%%%
%%%%%%%%
\renewcommand{\div}{\mbox{div }}
\newcommand{\curl}{\mbox{curl }}
\newcommand{\trchb}{\mbox{tr} \chib}
\def\trch{\mbox{tr}\chi}


\newcommand{\Ls}{{\mathcal L} \mkern-10mu /\, }
\newcommand{\eps}{{\epsilon} \mkern-8mu /\, }
%%%%%%%%%
%%%%%%%%%






\newcommand{\tr}{\mbox{tr}}


\newcommand{\xib}{\underline{\xi}}
\newcommand{\psib}{\underline{\psi}}
\newcommand{\rhob}{\underline{\rho}}
\newcommand{\thetab}{\underline{\theta}}
\newcommand{\gammab}{\underline{\gamma}}
\newcommand{\nub}{\underline{\nu}}
\newcommand{\lb}{\underline{l}}
\newcommand{\mub}{\underline{\mu}}
\newcommand{\Xib}{\underline{\Xi}}
\newcommand{\Thetab}{\underline{\Theta}}
\newcommand{\Lambdab}{\underline{\Lambda}}
\newcommand{\vphb}{\underline{\varphi}}

\newcommand{\ih}{\hat{i}}
\newcommand{\ui}{u_{\infty}}
\newcommand{\shb}{L^2_{sc}(\underline{H}_{\ub}^{(u_{\infty}, u)})}
%\newcommand{\sh}{L^2_{sc}(H_{u}^{(0, \ub)})}
\newcommand{\Rb}{\underline{\mathcal{R}}}
\newcommand{\tc}{\widetilde{\tr\chib}}

\newcommand{\tcL}{\widetilde{\mathscr{L}}}


\newcommand{\sRic}{Ric\mkern-19mu /\, \, \, \, }
\newcommand{\sL}{{\cal L}\mkern-10mu /}
\newcommand{\sLh}{\hat{\sL}}
\newcommand{\sg}{g\mkern-9mu /}
\newcommand{\seps}{\epsilon\mkern-8mu /}
\newcommand{\sd}{d\mkern-10mu /}
\newcommand{\sR}{R\mkern-10mu /}
\newcommand{\snab}{\nabla\mkern-13mu /}
\newcommand{\sdiv}{\mbox{div}\mkern-19mu /\, \, \, \, }
\newcommand{\scurl}{\mbox{curl}\mkern-19mu /\, \, \, \, }
\newcommand{\slap}{\mbox{$\triangle  \mkern-13mu / \, $}}
\newcommand{\scaletwoHuprime}[1]{\lVert #1 \rVert_{L^2_{(sc)}(H_{u^{\prime}}^{(0, \ubar)}    ) }  }
\newcommand{\sGamma}{\Gamma\mkern-10mu /}
\newcommand{\somega}{\omega\mkern-10mu /}
\newcommand{\somb}{\omb\mkern-10mu /}
\newcommand{\spi}{\pi\mkern-10mu /}
\newcommand{\sJ}{J\mkern-10mu /}
\renewcommand{\sp}{p\mkern-9mu /}
\newcommand{\su}{u\mkern-8mu /}
\newcommand{\m}{\mu}
\newcommand{\n}{\nu}
%%%%	
\newcommand\restri[2]{{% we make the whole thing an ordinary symbol
		\left. \kern-\nulldelimiterspace % automatically resize the bar with \right
		#1 % the function
		%\vphantom{\big|} % pretend it's a little taller at normal size
		\right|_{#2} % this is the delimiter
}}
%\newcommand{\paola}[1]{\textcolor{blue}{#1}}

\definecolor{ffqqqq}{rgb}{1. , 0. , 0. }
\definecolor{uuuuuu}{rgb}{0. 26666666666666666, 0. 26666666666666666, 0. 26666666666666666}

%\DeclareMathOperator{\argmin}{argmin}
\DeclareMathOperator{\sgn}{sgn}
\DeclareMathOperator{\dist}{dist}
\DeclareMathOperator{\diam}{diam}
%\DeclareMathOperator{\proj}{Proj}\includegraphics[]{. . /. . /. . /. . /Desktop/NiceMunich. pdf}


\delimitershortfall=-0. 1pt

\makeatletter
\def\ps@pprintTitle{%
	\let\@oddhead\@empty
	\let\@evenhead\@empty
	\let\@oddfoot\@empty
	\let\@evenfoot\@oddfoot
}
\def\@author#1{\g@addto@macro\elsauthors{\normalsize%
		\def\baselinestretch{1}%
		\upshape\authorsep#1\unskip\textsuperscript{%
			\ifx\@fnmark\@empty\else\unskip\sep\@fnmark\let\sep=, \fi
			\ifx\@corref\@empty\else\unskip\sep\@corref\let\sep=, \fi
		}%
		\def\authorsep{\unskip, \space}%
		\global\let\@fnmark\@empty
		\global\let\@corref\@empty  %% Added
		\global\let\sep\@empty}%
	\@eadauthor={#1}
}
\makeatother
\fi

%\newcommand{\norm}[1]{\left\Vert#1\right\Vert}
%\newcommand{\abs}[1]{\left\vert#1\right\vert}
%\newcommand{\set}[1]{\left\{#1\right\}}
%\newcommand{\real}{\mathbb R^n}
%\newcommand{\R}{\mathbb R}
%\newcommand{\eps}{\varepsilon}
%\newcommand{\To}{\longrightarrow}
\newcommand{\pf}[2]{\frac{\partial #1}{\partial #2}}
\newcommand{\wc}{\llcorner}
\newcommand{\cta}{\theta}
\newcommand{\Om}{\Omega}
\newcommand{\om}{\omega}
\newcommand{\vp}{\varphi}
\newcommand{\vr}{\varrho}
\newcommand{\sub}{\subset}
\newcommand{\half}{\frac{1}{2}}

%%%%%% My Own Commands
\newcommand{\tab}[1]{
	\begin{tabular}[]{||p{11cm}||}
		#1
\end{tabular}}

\newcommand{\foo}[1]{\footnote{#1}}
\newcommand{\red}[1]{ {\color{red} #1} }
\newcommand{\blue}[1]{ {\color{blue} #1} }
%\newcommand{\tr}{\mbox{\rm tr\, }}

\newcommand{\thmref}[1]{Theorem~\ref{#1}}
\newcommand{\secref}[1]{Section~\ref{#1}}
\newcommand{\subsecref}[1]{Subsection~\ref{#1}}
\newcommand{\lemref}[1]{Lemma~\ref{#1}}
\newcommand{\defref}[1]{Definition~\ref{#1}}
\newcommand{\propref}[1]{Proposition~\ref{#1}}
\newcommand{\probref}[1]{Problem~\ref{#1}}
\newcommand{\corref}[1]{Corollary~\ref{#1}}
\newcommand{\conjref}[1]{Conjecture~\ref{#1}}
\newcommand{\exampref}[1]{Example~\ref{#1}}
\newcommand{\quesref}[1]{Question~\ref{#1}}
\newcommand{\remref}[1]{Remark~\ref{#1}}

\iffalse
\usepackage{pdfpages}



%%%%%%%%%%%%%%%%%%%%%%%%%
%
% an alternative font
%
%  Latex-ing with this takes much longer
%
%\usepackage{fouriernc}
%\usepackage[T1]{fontenc}
%
%
%%%%%%%%%%%%%%%%%%%%%%%%%



%
%% Cyrillic letters
%
%  see page 20 in: https://www. latex-project. org/help/documentation/encguide. pdf
%
\DeclareFontFamily{U}{wncy}{}
\DeclareFontShape{U}{wncy}{m}{n}{<->wncyr10}{}
\DeclareSymbolFont{mcy}{U}{wncy}{m}{n}
\DeclareMathSymbol{\Sha}{\mathord}{mcy}{"58} % unused
\DeclareMathSymbol{\Chu}{\mathord}{mcy}{"51}
\DeclareMathSymbol{\Tsu}{\mathord}{mcy}{"43} % unused
\DeclareMathSymbol{\Yu}{\mathord}{mcy}{"10}
\DeclareMathSymbol{\yu}{\mathord}{mcy}{"18}  % unused
\DeclareMathSymbol{\Zhu}{\mathord}{mcy}{"11} % unused
\DeclareMathSymbol{\zhu}{\mathord}{mcy}{"19}
\DeclareMathSymbol{\Ya}{\mathord}{mcy}{"17}
\DeclareMathSymbol{\cyrl}{\mathord}{mcy}{"6C} % unused
\DeclareMathSymbol{\cyrL}{\mathord}{mcy}{"4C}






%%%%%%%%%%%%%%%%%%%%%%%%%
%
% (use later at the end - it takes much longer to LateX with this)
%
%% an alternative font
%\usepackage{fouriernc}
%\usepackage[T1]{fontenc}
%
%%%%%%%%%%%%%%%%%%%%%%%%%


%%%%%%%%%%%%%%%%%%%%%%%
%% KF macros 
%%%%%%%%%%%%%%%%%%%%%%%
\usepackage{color}
\newcommand{\kcom}[1]{{\color{red}{Kevin: #1}} }
\newcommand{\jcom}[1]{{\color{red}{James: #1}} }
\newcommand{\red}[1]{{\color{red}{#1}}}
\newcommand{\blue}[1]{{\color{blue}{#1}}}
\usepackage{tikz}


\DeclareMathAlphabet{\curly}{U}{rsfs}{m}{n}  %% curly font

\textheight=8. 5in
\textwidth=6. 5in
\oddsidemargin=0pt
\evensidemargin=0pt
\hoffset=0in

\newcommand{\ZZ}{{\mathbb Z}}
\newcommand{\QQ}{{\mathbb Q}}
\newcommand{\RR}{{\mathbb R}}
\newcommand{\CC}{{\mathbb C}}
\newcommand{\NN}{{\mathbb N}}
\newcommand{\VV}{{\mathbb V}}  
\newcommand{\LL}{{\mathbb L}}  
\newcommand{\er}{\mathrm{e}}

%%
%% probability,  expectation
%%
\newcommand{\PR}{\mathbb P}
\newcommand{\E}{\mathbb E}
%% fixes a serious LaTeX error in displaying \pmod,  especially in sums
\makeatletter
\renewcommand{\pmod}[1]{\allowbreak\mkern7mu({\operator@font mod}\, \, #1)}
\makeatother
%% sinlge line,  with or without a label
\newcommand{\be}{\begin{equation}}
	\newcommand{\ee}{\end{equation}}

\renewcommand{\a}{\ensuremath{\alpha}}
\renewcommand{\b}{\ensuremath{\beta}}
\newcommand{\del}{\ensuremath{\delta}}
\newcommand{\eps}{\ensuremath{\varepsilon}}
\newcommand{\g}{\ensuremath{\gamma}}
\renewcommand{\le}{\leqslant}
\renewcommand{\leq}{\leqslant}
\renewcommand{\ge}{\geqslant}
\renewcommand{\geq}{\geqslant}

\newcommand{\ssum}[1]{\sum_{\substack{#1}}}  %%% stacked sum
\newcommand{\sprod}[1]{\prod_{\substack{#1}}}  %% stacked product
\newcommand{\mint}[1]{\idotsint\limits_{\substack{#1}}} %% mult integral with conditions
\newcommand{\meas}{\mathrm{meas}}


\newcommand{\fl}[1]{{\ensuremath{\left\lfloor {#1} \right\rfloor}}}
\newcommand{\cl}[1]{{\ensuremath{\left\lceil #1 \right\rceil}}}
\newcommand{\order}{\asymp}      %% order of magnitude
\renewcommand{\(}{\left(}
\renewcommand{\)}{\right)}
\newcommand{\ds}{\displaystyle}
\newcommand{\pfrac}[2]{\left(\frac{#1}{#2}\right)}  %%% frac with paren
\newcommand{\li}{\mathrm{li}}
\newcommand{\emptyvec}{\varnothing}  %% empty vector

\newcommand{\balpha}{\ensuremath{\boldsymbol{\alpha}}}
\newcommand{\bbeta}{\ensuremath{\boldsymbol{\beta}}}
\newcommand{\brho}{\ensuremath{\boldsymbol{\rho}}}
\newcommand{\bgam}{\ensuremath{\boldsymbol{\gamma}}}
\newcommand{\bphi}{\ensuremath{\boldsymbol{\phi}}}
\newcommand{\bpsi}{\ensuremath{\boldsymbol{\psi}}}
\newcommand{\bxi}{\ensuremath{\boldsymbol{\xi}}}
\newcommand{\bomega}{\ensuremath{\boldsymbol{\omega}}}

\newcommand{\one}{\ensuremath{\mathbbm{1}}}  %% indicator function
\newcommand{\bone}{\ensuremath{\mathbf{1}}}  %% bold 1 - special vector function

%%  better subscript placement for tall subscripts (capital letters,  1,  2,  k,  etc)
\newcommand{\ssc}[2]{\ensuremath{{#1}_{#2}^{\phantom{2}}}}
\newcommand{\epszero}{\ssc{\eps}{0}}

%===>> CALLIGRAPHIC ALPHABET

\newcommand{\cA}{\ensuremath{\mathcal{A}}}
\newcommand{\cB}{\ensuremath{\mathcal{B}}}
\newcommand{\cC}{\ensuremath{\mathcal{C}}}
\newcommand{\cD}{\ensuremath{\mathcal{D}}}
\newcommand{\cE}{\ensuremath{\mathcal{E}}}
\newcommand{\cF}{\ensuremath{\mathcal{F}}}
\newcommand{\cG}{\ensuremath{\mathcal{G}}}
\newcommand{\cH}{\ensuremath{\mathcal{H}}}
\newcommand{\cI}{\ensuremath{\mathcal{I}}}
\newcommand{\cJ}{\ensuremath{\mathcal{J}}}
\newcommand{\cK}{\ensuremath{\mathcal{K}}}
\newcommand{\cL}{\ensuremath{\mathcal{L}}}
\newcommand{\cM}{\ensuremath{\mathcal{M}}}
\newcommand{\cN}{\ensuremath{\mathcal{N}}}
\newcommand{\cO}{\ensuremath{\mathcal{O}}}
\newcommand{\cP}{\ensuremath{\mathcal{P}}}
\newcommand{\cQ}{\ensuremath{\mathcal{Q}}}
\newcommand{\cR}{\ensuremath{\mathcal{R}}}
\newcommand{\cS}{\ensuremath{\mathcal{S}}}
\newcommand{\cU}{\ensuremath{\mathcal{U}}}
\newcommand{\cT}{\ensuremath{\mathcal{T}}}
\newcommand{\cV}{\ensuremath{\mathcal{V}}}
\newcommand{\cW}{\ensuremath{\mathcal{W}}}
\newcommand{\cX}{\ensuremath{\mathcal{X}}}
\newcommand{\cY}{\ensuremath{\mathcal{Y}}}
\newcommand{\cZ}{\ensuremath{\mathcal{Z}}}

%%%%%%%%%%%%%%%%%%%%%%%%%%%%%%%%%%%%%%%%%%%
%===>> SCRIPT ALPHABET

\newcommand{\sA}{\ensuremath{\mathscr{A}}}
\newcommand{\sB}{\ensuremath{\mathscr{B}}}
\newcommand{\sC}{\ensuremath{\mathscr{C}}}
\newcommand{\sD}{\ensuremath{\mathscr{D}}}
\newcommand{\sE}{\ensuremath{\mathscr{E}}}
\newcommand{\sF}{\ensuremath{\mathscr{F}}}
\newcommand{\sG}{\ensuremath{\mathscr{G}}}
\newcommand{\sH}{\ensuremath{\mathscr{H}}}
\newcommand{\sI}{\ensuremath{\mathscr{I}}}
\newcommand{\sJ}{\ensuremath{\mathscr{J}}}
\newcommand{\sK}{\ensuremath{\mathscr{K}}}
\newcommand{\sL}{\ensuremath{\mathscr{L}}}
\newcommand{\sM}{\ensuremath{\mathscr{M}}}
\newcommand{\sN}{\ensuremath{\mathscr{N}}}
\newcommand{\sO}{\ensuremath{\mathscr{O}}}
\newcommand{\sP}{\ensuremath{\mathscr{P}}}
\newcommand{\sQ}{\ensuremath{\mathscr{Q}}}
\newcommand{\sR}{\ensuremath{\mathscr{R}}}
\newcommand{\sS}{\ensuremath{\mathscr{S}}}
\newcommand{\sU}{\ensuremath{\mathscr{U}}}
\newcommand{\sT}{\ensuremath{\mathscr{T}}}
\newcommand{\sV}{\ensuremath{\mathscr{V}}}
\newcommand{\sW}{\ensuremath{\mathscr{W}}}
\newcommand{\sX}{\ensuremath{\mathscr{X}}}
\newcommand{\sY}{\ensuremath{\mathscr{Y}}}
\newcommand{\sZ}{\ensuremath{\mathscr{Z}}}
%% BOLD Letters

\newcommand{\ba}{\ensuremath{\mathbf{a}}}
\newcommand{\bb}{\ensuremath{\mathbf{b}}}
\newcommand{\bc}{\ensuremath{\mathbf{c}}}
\newcommand{\bC}{\ensuremath{\mathbf{C}}}
\newcommand{\bd}{\ensuremath{\mathbf{d}}}
\newcommand{\bee}{\ensuremath{\mathbf{e}}}
%\newcommand{\be}{\ensuremath{\mathbf{e}}}
%\newcommand{\bf}{\ensuremath{\mathbf{f}}}
\newcommand{\bg}{\ensuremath{\mathbf{g}}}
\newcommand{\bh}{\ensuremath{\mathbf{h}}}
\newcommand{\bi}{\ensuremath{\mathbf{i}}}
\newcommand{\bj}{\ensuremath{\mathbf{j}}}
\newcommand{\bk}{\ensuremath{\mathbf{k}}}
\newcommand{\bl}{\ensuremath{\mathbf{l}}}
\newcommand{\bmm}{\ensuremath{\mathbf{m}}}
\newcommand{\bM}{\ensuremath{\mathbf{M}}}
\newcommand{\bN}{\ensuremath{\mathbf{N}}}
\newcommand{\bn}{\ensuremath{\mathbf{n}}}
\newcommand{\bo}{\ensuremath{\mathbf{o}}}
\newcommand{\bp}{\ensuremath{\mathbf{p}}}
\newcommand{\bq}{\ensuremath{\mathbf{q}}}
\newcommand{\br}{\ensuremath{\mathbf{r}}}
\newcommand{\bs}{\ensuremath{\mathbf{s}}}
\newcommand{\bt}{\ensuremath{\mathbf{t}}}
\newcommand{\bu}{\ensuremath{\mathbf{u}}}
\newcommand{\bv}{\ensuremath{\mathbf{v}}}
\newcommand{\bw}{\ensuremath{\mathbf{w}}}
\newcommand{\bx}{\ensuremath{\mathbf{x}}}
\newcommand{\by}{\ensuremath{\mathbf{y}}}
\newcommand{\bz}{\ensuremath{\mathbf{z}}}


\newcommand{\tf}{\widetilde{f}}
\newcommand{\tl}{\ensuremath{\widetilde{\lambda}}}

\newcommand{\ft}{\tf}  %%% I write it both ways
%% multiple integral with multiple conditions

\newcommand{\scsc}[1]{\scriptscriptstyle{#1}}

%% Type II range
\newcommand{\two}{\ensuremath{[\theta, \theta+\nu]}}

%% C_{bounded}
\newcommand{\CB}{C_{\text{bd}}}
\fi
%\newcommand{\abs}[1]{\left\vert#1\right\vert}
%\newcommand{\norm}[1]{\Vert#1\Vert}
%\newcommand\F{\mathbb{F}}
\newcommand\E{\mathbb{E}}
%\newcommand\R{\mathbb{R}}
%\newcommand\Z{\mathbb{Z}}
%\newcommand\N{\mathbb{N}}
%\newcommand\n{\mathbf{n}}
%\newcommand\one{\mathbf{1}}
%\newcommand\C{\mathbb{C}}
%\newcommand\Q{\mathbb{Q}}
%\newcommand\eps{\varepsilon}
%\newcommand\B{\operatorname{B}}
%\newcommand\SO{{\operatorname{SO}}}
%\newcommand\Rot{{\operatorname{Rot}}}
%\newcommand\Dil{{\operatorname{Dil}}}
%\newcommand\df{{\operatorname{df}}}
\newcommand\Dcal{\mathcal{D}}
%\newcommand\DHL{\operatorname{DHL}}
%\newcommand\EH{\operatorname{EH}}
%\newcommand\MPZ{\operatorname{MPZ}}
\newcommand\TypeI{\operatorname{Type}_{\operatorname{I}}}
\newcommand\TypeII{\operatorname{Type}_{\operatorname{II}}}
\newcommand\TypeIII{\operatorname{Type}_{\operatorname{III}}}
\renewcommand\P{\mathbb{P}}

%\newcommand\rk{\mathrm{rk}}
%\newcommand\Fr{\mathrm{Fr}}
%\newcommand\tr{\mathrm{tr}}

\newcommand\ov{\overline}
\newcommand\wcheck{\widecheck}
%\newcommand\cond{\mathrm{cond}}
%\newcommand\Gal{\mathrm{Gal}}
%\newcommand\Kl{\mathrm{Kl}}
%\newcommand\swan{\mathrm{swan}}
%\newcommand\FT{\mathrm{FT}}
\newcommand\mcF{\mathcal{F}}
\newcommand\mcG{\mathcal{G}}
\newcommand\mcK{\mathcal{K}}
\newcommand\mcC{\mathcal{C}}
\newcommand\mcL{\mathcal{L}}
\newcommand\Gm{\mathbb{G}_m}
\newcommand\Af{\mathbb{A}^1}
\newcommand\Aa{\mathbb{A}}
\newcommand\Fp{\F_{p}}
\newcommand\Fq{\F_{q}}
\def \Re {\, {\rm Re}\, }
\def \supp {\, {\rm supp}\, }
\def\stacksum#1#2{{\stackrel{{\scriptstyle #1}}
		{{\scriptstyle #2}}}}
\newcommand{\mods}[1]{\, (\mathrm{mod}\, {#1})}
\newcommand{\piar}{\pi_1}
\newcommand{\pige}{\pi_1^{g}}
\newcommand\sumsum{\mathop{\sum\sum}\limits}
%\newcommand\E{\mathbb{E}}
\newcommand\Rr{\mathbb{R}}
\newcommand\Zz{\mathbb{Z}}
\newcommand\Nn{\mathbb{N}}
\newcommand\Cc{\mathbb{C}}
\newcommand\Qq{\mathbb{Q}}
\newcommand{\what}{\widehat}
%\newcommand{\ra}{\rightarrow}

\newcommand{\alert}[1]{{\bf \color{red} TODO: #1}}

\def\d{\mathrm{d}}


%\newcommand{\pian}[2]{\frac{\partial #1 }{\partial #2 }}
%\newcommand{\con}[2]{\tilde{\Gamma}_{#1}^{#2}}
\newcommand{\conn}[1]{\widetilde{\nabla}_{#1}}
\newcommand{\ykuo}[1]{\left( #1\right) }
\newcommand{\jkuo}[1]{\left\langle  #1\right\rangle }
%\newcommand{\fanshu}[1]{\left\|   #1\right\|}
%\newcommand{\covt}[1]{\tilde{\nabla}_{#1}}
%\newcommand{\F}{F^{\prime}\yuankuo{\frac{|du_{s,t}|^2}{2}}}
%\newcommand{\FF}{F^{\prime\prime}\yuankuo{\frac{|du_{s,t}|^2}{2}}}
%\newcommand{\pl}[1]{\Phi_*e_{#1}}
%\newcommand{\cov}[2]{\Phi_{*}\nabla_{e_{#1}}  e_{#2}  } 


%\newcommand{\curtensor}[1]{R^N\left(#1\right) }
%\newcommand{\cur}[2]{R_{#1}^{#2} }
\newcommand{\cur}[4]{#1_{#2}^{#3}\left( #4 \right) }

%\newcommand{\curva}[4]{R^{#1}_{#2,#3, #4} }
%\newcommand{\suanzi}[1]{\operatorname{#1}}
%\newcommand{\uu}[2]{u_{#1}^{#2}}
%\newcommand{\hessian}[2]{\nabla ^2 u (#1,#2)}
%\newcommand{\he}[2]{\sum_{#1}^{#2}}
%\newcommand{\jifen}[1]{\int_M #1 \mathrm{d}v_{g}}
%
%\newcommand{\n}{\nonumber}



%\usepackage{showkeys}


%自己的常用快捷输入
\def\s{\,\,\,\,}
\def\lan{\langle}
\def\ran{\rangle}
\def\dis{\displaystyle}
\def\dint{\displaystyle{\int}}
\def\mv{1.7ex}
\def\endproof{$\hfill\Box$\\}
\def\R{\mathbb{R}}
\def\H{\mathcal{H}}
\def\K{\mathcal{K}}
\def\C{\mathbb{C}}
%\def\vec#1{\vec{#1}}
\def\fu{f \circ u}
\def\dfu {\d f (du)}
\def \d {\mathrm{d}}
\def \Ric {\mathrm{Ric}}
\def \vol{\mathrm{Vol}}
\def \div{\mathrm{div}}
\def \hess{\mathrm{Hess}}
\def \S{\mathbb{S}}
\def \S{\mathbb{R}}
\def \o{\mathcal{O}}
\def\vecc#1{\vec{#1}}
\def\p{\partial}
\def\G{\widetilde{\Gamma}}
%\def\con#1#2{\widetilde{\Gamma}_{#1}^{#2}}
\def\g#1{\widetilde{g}_{#1}}
\def\a{\alpha}
%\def\b{\beta}wwwwwwwwwwwwwwwwwwwwwwwww
\def\ga{\gamma}
\def\e{\eta}
\def\v{V_\lambda^H}
\def\la{\lambda}
\def\si{\sigma}
\def\dfu#1{\d f (\d u (#1))}
\def\ep{\epsilon}
\def\dv{\d v_g}
%\defwwwwww
\usepackage{slashed}\usepackage[utf8]{inputenc}
%\usepackage[plainpages=false,pdfpagelabels=true]{hyperref}
\title{  Liouville theorem on  $p$-biharmonic map from gradient Ricci soliton}

\author{Xiang-Zhi Cao\thanks{School of Information Engineering, Nanjing Xiaozhuang University, Nanjing 211171, China}}

\numberwithin{equation}{section}


\usepackage{setspace}
\onehalfspacing
\setlength{\parskip}{0pt}
\setlength{\parindent}{2em}
\begin{document}
	\maketitle 
	\section{Introduction}
	
	Let $u:(M,g) \to (N,h)$ be  a smooth map  between Riemannian manifold. In the past decades, two famous map in geometry  has been studied widely, one is harmonic map(cf.\cite{MR164306}) , the other one is biharmonic map.  Jiang \cite{zbMATH05779520} defined biharmonic map,firstly. later, people began to generalize the concept of biharmonic map,such as  $F$-biharmonic map (cf.\cite{zbMATH06287169,zbMATH07642828} ), $f$-biharmonic maps and $f$-harmonic map, for example (cf. \cite{zbMATH06834772,zbMATH07076904,zbMATH06834772,zbMATH06149656,zbMATH05791335,zbMATH08101530,zbMATH06182078,zbMATH06149656}). $p$-biharmonic map (cf. \cite{MR4169439}) which is defined as $\tau_p(u)=\tau_p(u)=\div_g( \norm{\d u}^{p-2}\d u)=0$ . The critical point of 
	the functional 
	$$ E(u)=\int_M \frac{1}{2} \norm{\tau_p(u)}^2\d V $$  
	is called $p$-biharmonic map.
	
		
	Harmonic map is also generalized  widely. For example	$F$-harmonic map(cf. \cite{MR1820701,MR1700595}), $ f$-harmonic map (cf.\cite{MR3255005,MR3681112}).	
	Cherif\cite{zbMATH06527253}   defined $L$-harmonic map.  let $L:M\times N\times\R \to(0,\infty)$
	\begin{equation}
		E_L(\phi;D) = \int_D L\!\bigl(x,\phi(x),e(\phi)(x)\bigr)\, v_g,
	\end{equation}Its critical point (\cite{zbMATH06527253}) is defined as 
	\begin{equation}\label{eq:tauL}
		\tau_L(\phi) = L'_\phi\,\tau(\phi) + d\phi\!\bigl(\grad^M L'_\phi\bigr) - (\grad^N L)\circ\phi=0
	\end{equation}
	
In order to unify  many different kinds of harmonic and biharmonic map and inspired by action principle in unifying field theory,we introduce the following functional:
	\begin{equation}
		\begin{split}
			\int_DB\bigl(x,\phi_t(x),e(\phi_t)(x),\frac{\|\tau_L(\phi)\|^2}{2}\bigr)\,\d v_{g}
		\end{split}
	\end{equation}
	where, the function $B$ is defined by $B:M\times N\times\R  \times\R\to(0,\infty)$, $(x,y,r,s)\mapsto B(x,y,r,s)$.

Consider 

	\begin{equation}
		\begin{split}
			E_{p,q}(\phi)= \frac{1}{q}	\int_M \|\tau_p(\phi)\|^q \,\d v_{g}
		\end{split}
	\end{equation}	
Its critical point is called  $p$-$q$-biharmonic map. in particarly, $p$-$2$-biharmonic map is just the $p$-biharmonic mentioned above. 	Obviously, $p$-$q$-biharmonic map is the special case  $B$-$L$ map
	

	As we know , Ricci soliton are characterized by the following equation: 
	\begin{equation}
		\label{eq:soliton}
		\operatorname{Ric^M}+L_Xg=\lambda g
	\end{equation}
	Here, $\operatorname{Ric^M}$ represents the Ricci curvature of the Riemannian 
	manifold $(M,g)$,  $ \lambda\in\R$. 	Ricci solitons  with $X=\na f $ are  \emph{gradient Ricci solitons}.
	
	\begin{lem}[cf.\cite{MR2448435}]\label{17}
		Assume that $(M,g,f)$ is a gradient Ricci soliton.
		Then the following equation holds
		\begin{equation}
			\label{identity-soliton-const}
			\operatorname{Scal}^M+\|\nabla f\|^2-2\lambda f=C
		\end{equation}
		for some constant C.
	\end{lem}
	
	
	\begin{lem}[cf. \cite{MR2448435}]\label{16}
		\label{lem:bound-nabla-f}
		Assume that (M,g,f) is a steady gradient Ricci soliton, then
		$	\operatorname{Scal}^M\geq 0 $ and 
		\begin{equation}
			\label{eq:steady-bound-f}
			\|\nabla f\|^2\leq C
		\end{equation}
		for some positive constant C.
	\end{lem}
	
	Branding \cite{MR4674270} studied Liouville theorem of harmonic map and biharmonic map from gradient Ricci soliton. Later, \cite{MR4993852} generalized it to the case of p-harmonic map from gradient Ricci soliton.
	
	
	 Inspired by these works, we want to generalize it  to the case of $p$-biharmonic map. By Theorem \ref{18}, we can get 	
	\begin{thm}
		Assume that $(M,g,f)$ is an gradient Ricci soliton  and
		let $\phi\colon M\to N$ be a smooth $4$-biharmonic map  with $\phi\in W^{2,4}(M , N) \cap W^{2,14}(M , N)$.		If 
		\begin{equation}
			\begin{split}			
				&4\int_M\eta^2\operatorname{Ric}^M(e_i,e_j)\langle\na \left(\norm{d\phi}^{2} d\phi \right)(e_i,e_j),\tau_4(\phi)\rangle \d v_g \\
				&+2\int_M \left\langle \d \phi,\na \tau_4(\phi) \right\rangle\left\langle \d \phi(e_i),\d \phi(e_j) \right\rangle\Ric_{ij} \d v_g \geq 0.\\
			\end{split}
		\end{equation}
		and
		\begin{equation}
			\begin{split}
				\operatorname{Scal}^M<\lambda(m-4)
			\end{split}
		\end{equation}
		Then, $u$ is $p$-harmonic map.
	\end{thm}
	
	
	\begin{thm}\label{17}
		Assume that $(M, g, f)$ is the two-dimensional cigar soliton defined by $(\mathbb{R}^2, \frac{\d x^2+\d y^2}{1+x^2+y^2},-\log(1+x^2+y^2))$.
		Let $\phi: M \to N$ be a smooth $4$-biharmonic map with 
			\begin{equation}\label{18}
			\int_M(\|d\phi\|^{10}+\|\na d\phi\|^2)\dv+\int_M\norm{\d \phi}^{4} \|\na d\phi\|^2 \dv<\infty.
		\end{equation} Then $\phi$ is $4$-harmonic.
	\end{thm}
	\begin{rem}
		One can obtain similar  results to Theorem 1.1 and Theorem 1.2 for B-L-maps and p-biharmonic map. However the conditons may be very complicated. It is left to the readers as an exercise.
	\end{rem}
	
	
	{\bf Notations:} We denote the scalar curvature of $(M,g)$  by $\Scal^M$. We write $\tau_p(\phi)=\div_g(\norm{\d \phi}^{p-2}\d \phi)$.


In section 2,we studied  the stress-energy tensor of $B$-$L$ map and $p$-$q$-biharmonic map. In secton 3, we studied $4$-biharmonic map.
\section{$B$-$L$-map and $p $-$q$-biharmonic map}






\subsection{The first variation formula}

$B^\prime_\phi,\, B^{\prime\prime}_\phi\in C^\infty(M)$ represent the first order partial derivative and the second order   partial derivative of $B$ with respct to $e(u)$. The first order partial derivative and the second order   partial derivative $B$ in terms of  $\frac{\|\tau_L(\phi)\|^2}{2}$  are written as  $
B^\prime_\tau(x)$,$B^{\prime\prime}_\tau(x)$, respectively.


By the starndard method, we can get


\begin{thm} Let $\phi_t$ be the smooth variation,$\phi_0=\phi$, 
	$v = \left.\frac{\partial\phi_t}{\partial t}\right|_{t=0}$ , then
	
	\begin{equation}
		\begin{split}
			\left.\frac{d}{dt}E_{B}(\phi_t;D)\right|_{t=0}	=-\int_M \left\langle \tau_B,v \right\rangle-\int_D \left\langle \tau_{2,L,B} ,v \right\rangle
		\end{split}
	\end{equation}
	
	where  \begin{equation}
		\begin{split}
			\tau_{2,L,B}(\phi)
			&= -L'_\phi\,\tr R^N(B^\prime_\tau\tau_L(\phi),d\phi)d\phi
			- \tr\nabla^\phi L'_\phi\nabla^\phi\left(  B^\prime_\tau\tau_L(\phi)  \right)\\
			&\quad + \bigl(\nabla^N_{B^\prime_\tau\tau_L(\phi)}\grad^N L\bigr)\circ\phi
			+ \langle\nabla^\phi \left( B^\prime_\tau\tau_L(\phi) \right),d\phi\rangle\,(\grad^N L')\circ\phi \\
			&\quad - \tr\nabla^\phi\langle\nabla^\phi\left(B^\prime_\tau \tau_L(\phi) \right),d\phi\rangle\,L''_\phi\,d\phi.
		\end{split}
	\end{equation}
	and
	\begin{equation}
		\begin{split}
			\tau_B=	B'_\phi\,\tau(\phi) + d\phi\!\bigl(\grad^M B'_\phi\bigr) - (\grad^N B)\circ\phi
		\end{split}
	\end{equation}	
	
	
\end{thm}

\begin{proof}
By\cite{zbMATH06527253}, we derive 
	\begin{equation}
		\begin{split}
			\left.\frac{d}{dt}E_{B}(\phi_t;D)\right|_{t=0}
			=& \int_D \left.\partial_t\!\Bigl(B\bigl(x,\phi_t(x),e(\phi_t)(x),\tau_{L}(\phi)\bigr)\Bigr)\right|_{t=0}v_g.\\
			=&-\int_D h(\tau_L(\phi),v)\,v_g-\int_DB^\prime_\tau h(\nabla^\varphi_{\partial_t}\tau_L(\phi_t),\tau_L(\phi_t))\,v_g.
		\end{split}
	\end{equation}
	
by the reasoning in the proof of biharmonic map, we know that 	
	\begin{equation}
		\begin{split}
			\int_DB^\prime_\tau h(\nabla^\varphi_{\partial_t}\tau_L(\phi_t),\tau_L(\phi_t))\,v_g=\int_D \left\langle \tau_{2,L,B} ,v \right\rangle
		\end{split}
	\end{equation}
\end{proof}


Let $g_t$ be the smooth variation of $(M,g)$.

\begin{equation}
	\begin{split}
		\left.\frac{d}{dt}E_B(\phi;D)\right|_{t=0}
		= \int_D \delta\!B\bigl(x,\phi_t(x),e(\phi_t)(x),\frac{\|\tau_L(\phi)\|^2}{2}\bigr)\,v_g
		\\
		+ \int_DB\bigl(x,\phi_t(x),e(\phi_t)(x),\frac{\|\tau_L(\phi)\|^2}{2}\bigr)\,\delta(v_{g_t}).
	\end{split}
\end{equation}
So,
\begin{equation}
	\begin{split}
		&\left.\frac{d}{dt}E_B(\phi;D)\right|_{t=0}\\
		=& \int_D \delta(e(\phi))\,B'_\phi\,v_g
		+\frac{1}{2}\int_D B^\prime_\tau\delta(\|\tau_L(\phi)\|^2)\,v_g	+ \int_D B\,\delta(v_{g_t}),\\
	\end{split}
\end{equation}
Let $\langle\,,\,\rangle$ denote the induced Riemannian metric on $\otimes^2 T^*M$. We have
\begin{equation}\label{eq:delta}
	\delta(e(\phi)) = -\frac{1}{2}\langle\phi^*h,\delta g\rangle,\qquad
	\delta(v_{g_t}) = \frac{1}{2}\langle g,\delta g\rangle\,v_g,
\end{equation}
where $\phi^*h$ is the pull-back of the metric $h$.

\begin{thm}\label{thm:stressL} Let $g_t$ be the smooth variation of $(M,g)$.
	\begin{equation}
		\begin{split}
			&\left.\frac{d}{dt}E_L(\phi;D)\right|_{t=0}\\
			=& \frac{1}{2}\int_D B'_\phi \langle L_\phi\,g - L'_\phi\,\phi^*h,\,\delta g\rangle\, \d v_g+\frac{1}{2} \int_D B^\prime_\tau \left\langle S_{2,L}(\phi) ,\delta g \right\rangle \d v_g
		\end{split}
	\end{equation}
	where $S_{2,L} $ is given by (cf.\cite{zbMATH06527253})
	
	\begin{equation}\label{eq:S2L}
		\begin{split}
			S_{2,L}(\phi)(X,Y)
			&= -\frac{1}{2}|\tau_L(\phi)|^2\,g(X,Y)
			- L'_\phi\langle d\phi,\nabla^\phi\tau_L(\phi)\rangle\,g(X,Y) \\
			&\quad + L'_\phi\,h(d\phi(X),\nabla^\phi_Y\tau_L(\phi))
			+ L'_\phi\,h(d\phi(Y),\nabla^\phi_X\tau_L(\phi)) \\
			&\quad - \tau_L(\phi)(L)\,g(X,Y)
			- \tau_L(\phi)(L')\,h(d\phi(X),d\phi(Y)) \\
			&\quad + L''_\phi\langle d\phi,\nabla^\phi\tau_L(\phi)\rangle\,h(d\phi(X),d\phi(Y)).
		\end{split}
	\end{equation}
\end{thm}

So we can get the stress energy tensor $S_{B,L}$  as 
\begin{equation}
	\begin{split}
		S_{B,L}	=B'_\phi \left(  L_\phi\,g - L'_\phi\,\phi^*h \right)+B^\prime_\tau  S_{2,L}(\phi) 
	\end{split}
\end{equation}




\subsection{$p$-$q$-biharmonic map}



\begin{thm}Let $\phi_t$ be the smooth variation,$\phi_0=\phi$, 
	$v = \left.\frac{\partial\phi_t}{\partial t}\right|_{t=0}$ , then
	\begin{equation}
		\begin{split}
			\frac{d}{d t} E_{p,q}(\phi_t)|_{t=0}=-\int_M \left\langle \tau_{2,p,q}, v \right\rangle
		\end{split}
	\end{equation}
	where
	\begin{equation}
		\begin{split}
			\tau_{2,p,q}=&-\norm{\d \phi}^{p-2}\tr_gR^N(\norm{\tau_p}^{q-2}\tau_p,\d \phi)\d \phi-\tr_g   \na\left[ \norm{\d \phi}^{p-2} \na \left( \norm{\tau_p}^{q-2}\tau_p(\phi) \right)\right] \\
			&-(p-2)\tr_g\left[ \na \left\langle \na \left(\norm{\tau_p}^{q-2} \tau_p(\phi) \right), \d \phi \right\rangle \norm{\d \phi}^{p-4}\d \phi\right] 
		\end{split}
	\end{equation}
	
	In the case q=2, we know that 
	\begin{equation}
		\begin{split}
			\tau_{2,p,2}=&-\norm{\d \phi}^{p-2}\tr_gR^N(\tau_p,\d \phi)\d \phi-\tr_g \na \norm{\d \phi}^{p-2}\na \tau_p(\phi)\\
			&-(p-2)\tr_g\na \left\langle \na \tau_p(\phi), \d \phi \right\rangle \norm{\d \phi}^{p-4}\d \phi
		\end{split}
	\end{equation}
\end{thm}

Now we consider the stress energy tensor.
\begin{lem}[cf.\cite{zbMATH06268870}]
	\begin{equation}
		\begin{split}
			&\int_D \norm{\d \phi}^{p-2}\left\langle \d \phi(\xi), \tau_p(\phi) \right\rangle\\
			=&\int_D\left\langle -\text{sym}(\na \norm{\d \phi}^{p-2}\left\langle \d \phi, \tau_p(\phi) \right\rangle) , \delta g\right\rangle \\
			+&\frac{1}{2} \int_D \left\langle \div^M\left( \norm{\d \phi}^{p-2}\left\langle \d \phi(\xi), \tau_p(\phi) \right\rangle^{\sharp}  \right)g, \delta g \right\rangle
		\end{split}
	\end{equation}
\end{lem}

\begin{lem}[cf.\cite{zbMATH06268870}]
	let $\omega=	\norm{\d \phi}^{p-4} \left\langle \d \phi,\tau_p \right\rangle$
	\begin{equation}
		\begin{split}
			-\int_D	\norm{\d \phi}^{p-4} \left\langle  \d \phi (\na \left\langle  \phi^*h,\delta g \right\rangle  ),\tau_p\right\rangle 
			=\int_D \left\langle \phi^*h,\delta g \right\rangle \div \omega
		\end{split}
	\end{equation}
\end{lem}




\begin{thm}Let $g_t$ be the smooth variation of $(M,g)$.
	\begin{equation}
		\begin{split}
			\frac{\d}{\d t}	E_{p,q}(\phi)=\frac{1}{2}\int_D \left\langle S_{2,p,q} , \delta g \right\rangle\d v_g
		\end{split}
	\end{equation}
	where \begin{equation}
		\begin{split}
			S_{2,p,q}(X,Y)=&\norm{\tau_p}^{q-2}\bigg[ -\big(\frac{1}{2}\|\tau_p(\phi)\|^{2}+\norm{\d \phi}^{p-2}\langle d\phi,\na\tau_p(\phi)\rangle\big)g(X,Y)\\
			&+\norm{\d \phi}^{p-2}\langle d\phi(X),\na_Y\tau_p(\phi)\rangle+\norm{\d \phi}^{p-2}\langle d\phi(Y),\na_X\tau_p(\phi)\rangle\\
			&+(p-2)\norm{\tau_p(\phi)}^{p-4}\left\langle \d \phi,\na \tau_p(\phi) \right\rangle\left\langle \d \phi(X),\d \phi(Y) \right\rangle\bigg]
		\end{split}
	\end{equation}
\end{thm}



\begin{proof}By standard process, we can show 
	\begin{equation}
		\begin{split}
			\left.\frac{d}{dt}E_{2,p,q}(\phi;D)\right|_{t=0}
			=& \int_D \delta  \frac{1}{q}	 \|\tau_p(\phi)\|^q\d v
			+ \int_D \frac{1}{q}	 \|\tau_p(\phi)\|^q,\delta(v_{g_t})\\
			=& \int_D   \frac{1}{2} \|\tau_p(\phi)\|^{q-2}	 \delta\|\tau_p(\phi)\|^2\d v
			+ \int_D \frac{1}{q}	 \|\tau_p(\phi)\|^q,\delta(v_{g_t})..
		\end{split}
	\end{equation}
	
Noticing (cf. Cherif \cite{zbMATH07894526})  has shown that 
	
	\begin{equation}
		\begin{split}
			\delta\|\tau_p(\phi)\|^2=&-(p-2)\norm{\d \phi}^{p-4}\left\langle  \phi^*h,\delta g \right\rangle \left\langle \tau(\phi), \tau_p(\phi)\right\rangle \\
			&-2\norm{\d \phi}^{p-2}\left\langle h( \na \d \phi,\tau_p(\phi)) , \delta g \right\rangle -2\norm{\d \phi}^{p-2}\left\langle \d \phi(\xi), \tau_p(\phi) \right\rangle \\
			&-(p-2)(p-4)\norm{\d \phi}^{p-5}\left\langle  \phi^*h,\delta g \right\rangle \left\langle \d \phi (\na |\d \phi|) , \tau_p(\phi) \right\rangle \\
			&-2(p-2)\norm{\d \phi}^{p-3} \left\langle \d |\d \phi|\otimes h(\d \phi, \tau_p(\phi)) ,\delta g \right\rangle \\
			&-(p-2)\norm{\d \phi}^{p-4} \left\langle  \d \phi (\na \left\langle  \phi^*h,\delta g \right\rangle  ),\tau_p\right\rangle 
		\end{split}
	\end{equation}
Following the proof of  Cherif\cite{zbMATH07894526}, we are done.
	
\end{proof}

	
\section{$p$-biharmonic map}	
 One can refer to (cf.[cf.\cite{zbMATH07894526}]) for the stress-energy tensor  associated with $p$-biharmonic map ,which is given by  
\begin{equation}
	\begin{split}
			S_{2,p}(X,Y)=&-\big(\frac{1}{2}\|\tau_p(\phi)\|^{2}+\norm{\d \phi}^{p-2}\langle d\phi,\na\tau_p(\phi)\rangle\big)g(X,Y)\\
		&+\norm{\d \phi}^{p-2}\langle d\phi(X),\na_Y\tau_p(\phi)\rangle+\norm{\d \phi}^{p-2}\langle d\phi(Y),\na_X\tau_p(\phi)\rangle\\
		&+(p-2)\norm{\tau_p(\phi)}^{p-4}\left\langle \d \phi,\na \tau_p(\phi) \right\rangle\left\langle \d \phi(X),\d \phi(Y) \right\rangle
	\end{split}
\end{equation}

\begin{lem}[cf.\cite{zbMATH07894526}]
	\begin{equation}
		\begin{split}
			\div_g S_{2,p}(X)=-\left\langle \tau_{2,p} , \d \phi(X)  \right\rangle
		\end{split}
	\end{equation}
\end{lem}
A routine computation gives 
	$$
\begin{aligned}
\Tr(S_{2,p})=&-m\big(\frac{1}{2}\|\tau_p(\phi)\|^{2}+\norm{\d \phi}^{p-2}\langle d\phi,\na\tau_p(\phi)\rangle\big)\\
&+\norm{\d \phi}^{p-2}\langle d\phi,\na\tau_p(\phi)\rangle+\norm{\d \phi}^{p-2}\langle d\phi,\na\tau_p(\phi)\rangle\\
&+(p-2)\norm{\tau_p(\phi)}^{p-4}\left\langle \d \phi,\na \tau_p(\phi) \right\rangle\norm{\na \phi}^2	
\end{aligned}
$$


 By adapting the proof  in \cite{MR4674270}  \cite{MR4993852}, it is not hard to get
\begin{lem}\label{im-1}[cf.\cite{zbMATH07894526}]
	Assume that $(M,g,f)$ is a gradient Ricci soliton and
	let $\phi\colon M\to N$ be a smooth biharmonic map.
	For $\eta\in C^\infty(M)$ compactly supported, then 
\begin{equation}\label{13}
	\begin{split}
			&\int_M\eta^2\big(\lambda(m-4)-\operatorname{Scal}^M\big)\|\tau_p(\phi)\|^2\dv\\
			&+4\int_M\eta^2\operatorname{Ric}^M(e_i,e_j)\langle\na \left(\norm{d\phi}^{p-2} d\phi \right)(e_i,e_j),\tau_p(\phi)\rangle \dv \\
		&+\int_M 2(p-2)\norm{\tau_p(\phi)}^{p-4}\left\langle \d \phi,\na \tau_p(\phi) \right\rangle\left\langle \d \phi(e_i),\d \phi(e_j) \right\rangle\Ric_{ij}\\&-\int_M \eta^2(p-2)\la\norm{\tau_p(\phi)}^{p-4}\left\langle \d \phi,\na \tau_p(\phi) \right\rangle\norm{\na \phi}^2 \dv\\
		=&
		-4\int_M\norm{\d \phi}^{p-2}\langle d\phi(\operatorname{Ric}^M(\nabla\eta^2)),\tau_p(\phi)\rangle \dv 
		+4\lambda\int_M\langle d\phi(\nabla\eta^2),\tau_p(\phi)\rangle \dv 
		\\
		&-\int_M\langle\nabla\eta^2,\nabla f\rangle\|\tau_p(\phi)\|^2 \dv +2\int_M(\Delta\eta^2)\norm{\d \phi}^{p-2}\langle d\phi(\nabla f),\tau_p(\phi)\rangle \dv 
		\\&+4\int_M(\nabla_{e_i}\eta^2)\nabla_{e_j}f\norm{\d \phi}^{p-2}\langle\na d\phi(e_i,e_j),\tau_p(\phi)\rangle \dv,
	\end{split}
\end{equation}
	where $\{e_i\},i=1,\ldots,m$ represents an orthonormal basis of $TM$.
\end{lem}

\begin{proof}
	As $\eta$ has compact support,
\begin{equation}\label{9}
	\begin{split}
			0=-\int_M\eta^2\langle\nabla f,\operatorname{div} S_{2,p}\rangle \dv
		=\int_M\eta^2\langle\nabla^2f,S_{2,p}\rangle \dv
		+\int_M S_{2,p}(\nabla f,\nabla\eta^2) \dv.
	\end{split}
\end{equation}
By the definiton of gradient Ricci soliton,  we deal with the  term $\int_M\eta^2\langle\nabla^2f,S_{2,p}\rangle \dv$
  by $\int_M \left\langle g,S_{2,p} \right\rangle$ and  $\int_M \left\langle \Ric^M,S_{2,p} \right\rangle$  respectively.
\begin{equation}\label{10}
	\begin{split}
		&\int_M\eta^2\langle g,S_{2,p}\rangle \dv\\
		=&\int_M\eta^2\tr_gS_{2,p} \dv \\
		=&\int_M\eta^2\bigg(-m\big(\frac{1}{2}\|\tau_p(\phi)\|^{2}+\norm{\d \phi}^{p-2}\langle d\phi,\na\tau_p(\phi)\rangle\big)\\
		&+\norm{\d \phi}^{p-2}\langle d\phi,\na\tau_p(\phi)\rangle+\norm{\d \phi}^{p-2}\langle d\phi,\na\tau_p(\phi)\rangle\\
		&+(p-2)\norm{\tau_p(\phi)}^{p-4}\left\langle \d \phi,\na \tau_p(\phi) \right\rangle\norm{\na \phi}^2\bigg) \dv\\
		=&\big(2-\frac{m}{2}\big)\int_M\eta^2\|\tau_p(\phi)\|^2 \dv
		+(2-m)\int_M\langle d\phi(\nabla\eta^2),\tau_p(\phi)\rangle \dv.\\
		&+\int_M \eta^2(p-2)\norm{\tau_p(\phi)}^{p-4}\left\langle \d \phi,\na \tau_p(\phi) \right\rangle\norm{\na \phi}^2\bigg) \dv
	\end{split}
\end{equation}
	Moreover, a similar calculation shows
\begin{equation}\label{11}
	\begin{split}
		&\int_M\eta^2\langle\operatorname{Ric}^M,S_{2,p}\rangle \dv\\
		=&\int_M-\eta^2\operatorname{Scal}^M \big(\frac{1}{2}\|\tau_p(\phi)\|^{2}+\norm{\d \phi}^{p-2}\langle d\phi,\na\tau_p(\phi)\rangle\big)\\
		&+\Ric_{ij}\norm{\d \phi}^{p-2}\langle d\phi(e_i),\na_{e_j}\tau_p(\phi)\rangle+\norm{\d \phi}^{p-2}\langle d\phi(e_j),\na_{e_i}\tau_p(\phi)\rangle\\
		&+(p-2)\norm{\tau_p(\phi)}^{p-4}\left\langle \d \phi,\na \tau_p(\phi) \right\rangle\left\langle \d \phi(e_i),\d \phi(e_j)\Ric_{ij} \right\rangle\\
		=&-\frac{1}{2}\int_M\eta^2\operatorname{Scal}^M\|\tau_p(\phi)\|^2 \dv
		+2\int_M\eta^2\operatorname{Ric}^M(e_i,e_j)\langle\na \left( \norm{\d \phi}^{p-2}d\phi(e_i,e_j) \right),\tau_p(\phi)\rangle \dv\\
		&+2\int_M\norm{\d \phi}^{p-2}\eta^2\langle d\phi\big(\operatorname{div}\operatorname{Ric}^M-\frac{1}{2}\nabla \operatorname{Scal}^M\big),\tau_p(\phi)\rangle \dv \\
		&-\int_M\norm{\d \phi}^{p-2} \operatorname{Scal}^M
		\langle d\phi(\nabla\eta^2),\tau_p(\phi)\rangle \dv
		+2\int_M\norm{\d \phi}^{p-2}\langle d\phi(\operatorname{Ric}^M(\nabla\eta^2)),\tau_p(\phi)\rangle \dv.\\
		&+\int_M (p-2)\norm{\tau_p(\phi)}^{p-4}\left\langle \d \phi,\na \tau_p(\phi) \right\rangle\left\langle \d \phi(e_i),\d \phi(e_j) \right\rangle\Ric_{ij}
	\end{split}
\end{equation}
	
Let $\{e_i\},i=1,\ldots,m$ be an orthonormal basis of $TM$ that satisfies 
	$\nabla_{e_i}e_j=0,i,j=1,\ldots,m$ at a fixed point $p\in M$.
	\begin{align*}
		&S_{2,p}(\nabla f,\nabla\eta^2)\\
		=-&\langle\nabla\eta^2,\nabla f\rangle\big(\frac{1}{2}\|\tau_p(\phi)\|^2
		+\norm{\d \phi}^{p-2}\langle d\phi,\na\tau_p(\phi)\rangle\big)\\
		&+\norm{\d \phi}^{p-2}(\nabla_{e_i}\eta^2)\nabla_{e_j}f\langle d\phi(e_i),\na_{e_j}\tau_p(\phi)\rangle \\
		&+\norm{\d \phi}^{p-2}(\nabla_{e_j}\eta^2)\nabla_{e_i}f\langle d\phi(e_i),\na_{e_j}\tau_p(\phi)\rangle.\\
		&+(p-2)\norm{\tau_p(\phi)}^{p-4}\left\langle \d \phi,\na \tau_p(\phi) \right\rangle\left\langle \d \phi(\nabla f),\d \phi(\nabla\eta^2) \right\rangle
	\end{align*}
	
	A direct calculation using integration by parts yields
\begin{equation}\label{5}
	\begin{split}
			&-\int_M\langle\nabla\eta^2,\nabla f\rangle\big(\frac{1}{2}\|\tau_p(\phi)\|^2
		+\norm{\d \phi}^{p-2}\langle d\phi,\na\tau_p(\phi)\rangle\big) \dv\\
		=&\frac{1}{2}\int_M\langle\nabla\eta^2,\nabla f\rangle\|\tau_p(\phi)\|^2 \dv \\
		&+\int_M(\nabla_{e_j}\nabla_{e_i}\eta^2)\nabla_{e_i}f\norm{\d \phi}^{p-2}\langle d\phi(e_j),\tau_p(\phi)\rangle \dv\\
		&+\int_M(\nabla_{e_i}\eta^2)\nabla_{e_j}\nabla_{e_i}f\norm{\d \phi}^{p-2}\langle d\phi(e_j),\tau_p(\phi)\rangle \dv.
	\end{split}
\end{equation}
	
	Moreover, using integration by parts once more we find
\begin{equation}\label{6}
	\begin{split}
		&\int_M(\nabla_{e_i}\eta^2)\nabla_{e_j} f\norm{\d \phi}^{p-2} \langle d\phi(e_i),\na_{e_j}\tau_p(\phi)\rangle \dv\\ 
		=&-\int_M(\nabla_{e_j}\nabla_{e_i}\eta^2)\nabla_{e_j} f\langle d\phi(e_i),\tau_p(\phi)\rangle \dv \\
		&-\int_M(\nabla_{e_i}\eta^2)\Delta f\langle d\phi(e_i),\tau_p(\phi)\rangle \dv \\
		&-\int_M(\nabla_{e_i}\eta^2)\nabla_{e_j} f\langle\na d\phi(e_i,e_j),\tau_p(\phi)\rangle \dv
	\end{split}
\end{equation}
	and 
\begin{equation}\label{7}
	\begin{split}
			&\int_M(\nabla_{e_j}\eta^2)\nabla_{e_i} f\norm{\d \phi}^{p-2}\langle d\phi(e_i),\na_{e_j}\tau_p(\phi)\rangle \dv \\
		=&-\int_M(\Delta\eta^2)\nabla_{e_i} f\langle d\phi(e_i),\tau_p(\phi)\rangle \dv \\
		&-\int_M(\nabla_{e_j}\eta^2)\nabla_{e_j}\nabla_{e_i}f
		\langle d\phi(e_i),\tau_p(\phi)\rangle \dv \\
		&-\int_M(\nabla_{e_j}\eta^2)\nabla_{e_i} f\langle\na(\norm{\d \phi}^{p-2}\nabla d\phi)(e_i,e_j),\tau_p(\phi)\rangle \dv.
	\end{split}
\end{equation}
	
	Adding up both contributions yields
\begin{equation}\label{8}
	\begin{split}
		&\int_M(\nabla_{e_i}\eta^2)\nabla_{e_j}f\langle d\phi(e_i),\na_{e_j}\tau_p(\phi)\rangle \dv 
		+\int_M(\nabla_{e_j}\eta^2)\nabla_{e_i}f\langle d\phi(e_i),\na_{e_j}\tau_p(\phi)\rangle \dv \\
		=&-\int_M(\nabla_{e_j}\nabla_{e_i}\eta^2)\nabla_{e_j}f\langle d\phi(e_i),\tau_p(\phi)\rangle \dv
		-\int_M(\nabla_{e_i}\eta^2)\Delta f\langle d\phi(e_i),\tau_p(\phi)\rangle \dv \\
		&-2\int_M(\nabla_{e_i}\eta^2)\nabla_{e_j}f\langle\na (\norm{\d \phi}^{p-2}\d\phi)(e_i,e_j),\tau_p(\phi)\rangle \dv\\
		&-\int_M(\Delta\eta^2)\nabla_{e_i}f\norm{\d \phi}^{p-2}\langle d\phi(e_i),\tau_p(\phi)\rangle \dv \\
		&-\int_M(\nabla_{e_j}\eta^2)\nabla_{e_j}\nabla_{e_i}f \norm{\d \phi}^{p-2} \langle d\phi(e_i),\tau_p(\phi)\rangle \dv.
	\end{split}
\end{equation}
	
By \eqref{5}\eqref{6}\eqref{7}\eqref{8}
\begin{equation}\label{12}
	\begin{split}
			&\int_M S_{2,p}(\nabla f,\nabla\eta^2)~\dv\\
		=&		-\frac{1}{2}\int_M\langle\nabla\eta^2,\nabla f\rangle\|\tau_p(\phi)\|^2 \dv 
		+\int_M(\Delta\eta^2)\langle d\phi(\nabla f),\tau_p(\phi)\rangle \dv \\
		\nonumber&+\int_M\Delta f\langle d\phi(\nabla\eta^2),\tau_p(\phi)\rangle \dv  
		+2\int_M(\nabla_{e_i}\eta^2)\nabla_{e_j}f\langle\na d\phi(e_i,e_j),\tau_p(\phi)\rangle \dv\\
		=&-\frac{1}{2}\int_M\langle\nabla\eta^2,\nabla f\rangle\|\tau_p(\phi)\|^2 \dv 
		+\int_M(\Delta\eta^2)\langle d\phi(\nabla f),\tau_p(\phi)\rangle \dv \\
		\nonumber&+\int_M(m\lambda-\operatorname{Scal}^M)\langle d\phi(\nabla\eta^2),\tau_p(\phi)\rangle \dv \\
		&+2\int_M(\nabla_{e_i}\eta^2)\nabla_{e_j}f\langle\na d\phi(e_i,e_j),\tau_p(\phi)\rangle \dv,
	\end{split}
\end{equation}
	where we used the equation for a gradient Ricci soliton \eqref{eq:soliton} 
	in the second step. The statements follows from \eqref{9}\eqref{10}\eqref{11}\eqref{12}.
\end{proof}

\begin{lem}\label{15}
	Assume that $(M,g,f)$ is a complete, non-compact gradient Ricci soliton
	with $\operatorname{Ric}^M\leq C$ and $\|\nabla f\|<\infty$.
	Let $\phi\colon M\to N$ be a smooth $p$-biharmonic map with 
	\begin{equation}\label{14}
		\int_M(\|d\phi\|^{4p-6}+\|\na d\phi\|^2)\dv+\int_M\norm{\d \phi}^{2p-4} \|\na d\phi\|^2 \dv<\infty.
	\end{equation}
	Then the following inequality holds
	\begin{equation}\label{20}
		\begin{split}
			&\int_M\big(\lambda(m-4)-\operatorname{Scal}^M\big)\|\tau_p(\phi)\|^2\dv
			+4\int_M\operatorname{Ric}^M(e_i,e_j)\langle\na d\phi(e_i,e_j),\tau_p(\phi)\rangle \dv\\
			&+\int_M (p-2)\norm{\tau_p(\phi)}^{p-4}\left\langle \d \phi,\na \tau_p(\phi) \right\rangle\left\langle \d \phi(e_i),\d \phi(e_j) \right\rangle\Ric_{ij}\\
			&+\int_M (p-2)\norm{\tau_p(\phi)}^{p-4}\left\langle \d \phi,\na \tau_p(\phi) \right\rangle\left\langle \d \phi(e_i),\d \phi(e_j) \right\rangle\Ric_{ij}\\&+\int_M \eta^2(p-2)\norm{\tau_p(\phi)}^{p-4}\left\langle \d \phi,\na \tau_p(\phi) \right\rangle\norm{\na \phi}^2\bigg) \dv\\
			&\leq 0.
		\end{split}
	\end{equation}
\end{lem}
\begin{proof}[Proof of Lemma \ref{15}]
	Again, let  $0\leq\eta\leq 1$ on $M$ be such that
	\begin{equation*}
		\eta(x)=1\textrm{ for } x\in B_R(x_0),\qquad \eta(x)=0\textrm{ for } x\in B_{2R}(x_0),\qquad |\nabla^q\eta|\leq\frac{C}{R^q}\textrm{ for } x\in M,
	\end{equation*}
	where $B_R(x_0)$ denotes the geodesic ball around the point $x_0$ with radius $R$ and $q=1,2$.
	Moreover, let $\{e_i\},i=1,\ldots,m$ be an orthonormal basis of $TM$ that satisfies 
	$\nabla_{e_i}e_j=0,i,j=1,\ldots,m$ at a fixed point $p\in M$.
	
	We need to estimate the terms on 	the right hand side of \eqref{13}. 	By the  assumption \eqref{14},	it is easy to infer
	\begin{align*}
		&\int_M \norm{\d \phi}^{p-2}\langle d\phi(\operatorname{Ric}^M(\nabla\eta^2)),\tau_p(\phi)\rangle \dv\leq\frac{C}{R}\int_M(\|\na d\phi\|^2+\|d\phi\|^{4p-6})\dv \to 0, \\
		&\int_M \norm{\d \phi}^{p-2} \langle d\phi(\nabla\eta^2),\tau_p(\phi)\rangle \dv 
		\leq\frac{C}{R}\int_M(\|\na d\phi\|^2+\|d\phi\|^{4p-6})\dv 
		\to 0,  \\
		&\int_M\langle\nabla\eta^2,\nabla f\rangle\|\tau_p(\phi)\|^2 \dv 
		\leq\frac{C}{R}\int_M \norm{\d \phi}^{2p-4}\|\na d\phi\|^2\dv 
		\to 0, \\
	&\int_M \norm{\d \phi}^{p-2}(\Delta\eta^2)\langle d\phi(\nabla f),\tau_p(\phi)\rangle \dv 
		\leq\frac{C}{R^2}\int_M(\|\na d\phi\|^2+\|d\phi\|^2)\dv 
		\to 0, \\
		&\int_M\norm{\d \phi}^{p-2}(\nabla_{e_i}\eta^2)\nabla_{e_j}f\langle\na d\phi(e_i,e_j),\tau_p(\phi)\rangle 
		\leq\frac{C}{R}\int_M\norm{\d \phi}^{2p-4} \|\na d\phi\|^2 \dv\to 0
	\end{align*}
	as $R\to\infty$.
	Note that we used the inequality $\|\tau_p(\phi)\|\leq C_p \norm{\d \phi}^{p-2}\|\na d\phi\|$ and $ \Ric\leq C, \norm{\na f}\leq C$.
	This completes the proof.
\end{proof}



By Lemma \ref{im-1}, we can get 
\begin{thm}\label{18}
	Assume that $(M,g,f)$ is an gradient Ricci soliton  and
	let $\phi\colon M\to N$ be a smooth $4$-biharmonic map with  condition 
	\begin{equation}\label{18}
		\int_M(\|d\phi\|^{10}+\|\na d\phi\|^2)\dv+\int_M\norm{\d \phi}^{4} \|\na d\phi\|^2 \dv+\int_M\norm{\d \phi}^{7} \|\na d\phi\|^2 \dv<\infty.
	\end{equation}
	 If 
	\begin{equation}
		\begin{split}			
			&4\int_M\eta^2\operatorname{Ric}^M(e_i,e_j)\langle\na \left(\norm{d\phi}^{2} d\phi \right)(e_i,e_j),\tau_4(\phi)\rangle \d v_g \\
			&+\int_M 2\left\langle \d \phi,\na \tau_4(\phi) \right\rangle\left\langle \d \phi(e_i),\d \phi(e_j) \right\rangle\Ric_{ij} \d v_g \geq 0.\\
		\end{split}
	\end{equation}
	and
	\begin{equation}
		\begin{split}
			\operatorname{Scal}^M<\lambda(m-4)
		\end{split}
	\end{equation}
Then, $u$ is $p$-harmonic map.
\end{thm}
\begin{proof}
We choose the same cutoff function $\eta$ as in the proof of Lemma \ref{15}. By \eqref{20}, we obtain that 
				\begin{equation}
				\begin{split}
						&\int_M\eta^2\big(\lambda(m-4)-\operatorname{Scal}^M\big)\|\tau_p(\phi)\|^2\dv+
					4\int_M\eta^2\operatorname{Ric}^M(e_i,e_j)\langle\na \left(\norm{d\phi}^{2} d\phi \right)(e_i,e_j),\tau_p(\phi)\rangle \dv \\
					&+\int_M 2\left\langle \d \phi,\na \tau_4(\phi) \right\rangle\left\langle \d \phi(e_i),\d \phi(e_j) \right\rangle\Ric_{ij}\\
					&+\int_M 2\eta^2\left\langle \norm{\na \phi}^2\d \phi,\na \tau_4(\phi) \right\rangle \dv\\
					\leq& 0.					
				\end{split}
			\end{equation}
	By intrgration by parts and the definiton of $\tau_4(\phi)$, we can show that 
		\begin{equation}
		\begin{split}
			&\int_M\eta^2\big(\lambda(m-4)-2\operatorname{Scal}^M\big)\|\tau_p(\phi)\|^2\dv\\
			&+
			4\int_M\eta^2\operatorname{Ric}^M(e_i,e_j)\langle\na \left(\norm{d\phi}^{2} d\phi \right)(e_i,e_j),\tau_p(\phi)\rangle \dv \\
									\leq  &\int_M 2\na\eta^2\left\langle \norm{\na \phi}^2\d \phi, \tau_4(\phi) \right\rangle \dv\\ 
										\end{split}
	\end{equation}
	Let $r\to \infty$, we find the right side tends to zero.	
\end{proof}

\begin{lem}[cf.\cite{MR954419}]\label{32}
	Let $(M,g)$ be steady gradient Ricci soliton, then
	\begin{equation}
		\begin{split}
			\norm{\na \Scal} \leq  \frac{C}{R}, \, \text{on}\, B_R(x_0)
		\end{split}
	\end{equation}
\end{lem}
\begin{proof}
	It follows from Lemma \ref{17},Lemma \ref{16}.
\end{proof}

\begin{proof}[Proof of Theorem \ref{17}]
	As we know cigar soliton is steady gradient Ricci soliton.
For two-dimensional cigar soliton	, we know that $\la=0,\Scal>0$. Obviously, $2\operatorname{Ric}^M=\operatorname{Scal}^Mg$.
	\begin{equation}
	\begin{split}
		&\int_M\eta^2\big(-2\operatorname{Scal}^M\big)\|\tau_p(\phi)\|^2\dv\\
		&+
		4\int_M\eta^2\operatorname{Ric}^M(e_i,e_j)\langle\na \left(\norm{d\phi}^{2} d\phi \right)(e_i,e_j),\tau_p(\phi)\rangle \dv \\
		\leq  &\int_M 2\na\eta^2\left\langle \norm{\na \phi}^2\d \phi, \tau_4(\phi) \right\rangle \dv\\ 
		+&\int_M 2\na\Scal^M\left\langle \norm{\na \phi}^2\d \phi, \tau_4(\phi) \right\rangle \dv
	\end{split}
\end{equation}

Lemma \ref{15} and Lemma \ref{16}, Lemma \ref{17} yields
\begin{equation*}
	\int_{R^2}\operatorname{Scal}^M\|\tau_p(\phi)\|^2\dv\leq 0.
\end{equation*}
As the cigar soliton has positive scalar curvature we can deduce $\tau_p(\phi)=0$
yielding the claim.
\end{proof}

	\bibliographystyle{plain}
	\bibliographystyle{plain}
\bibliography{a.bib,3-13.bib,zbmath-cherif.bib,mr3.15.bib,zbmath-3.15-2.bib,myonlymathscinetbibfrom2023.bib}
	
	
	\end{document}
